% =====================================================================
%  Mechanics III -- Lecture Notes
%  Written up from handwritten notes of Lectures 1--13 (Autumn 2026),
%  with full derivations and explanations.
% =====================================================================
\documentclass[11pt,a4paper,oneside]{report}

\usepackage[utf8]{inputenc}
\usepackage[T1]{fontenc}
\usepackage{lmodern}
\usepackage{microtype}
\usepackage[left=2.8cm,right=4.2cm,top=2.8cm,bottom=3cm,marginparwidth=2.7cm,marginparsep=0.5cm]{geometry}
\usepackage{amsmath,amssymb,amsthm,mathtools}
\usepackage{bm}
\usepackage{siunitx}
\usepackage{enumitem}
\usepackage{titlesec}
\usepackage{tikz}
\usetikzlibrary{arrows.meta,calc,patterns}
\usepackage[hidelinks]{hyperref}
\usepackage{cleveref}

\setlength{\parindent}{1.2em}
\setlength{\parskip}{0.25em}
\numberwithin{equation}{chapter}
\setlist{itemsep=0.15em,topsep=0.3em}

\titleformat{\chapter}[display]{\normalfont\bfseries}{\large\mdseries\scshape Chapter \thechapter}{6pt}{\huge}
\titlespacing*{\chapter}{0pt}{-10pt}{28pt}

% ---------- notation ----------
\newcommand{\vect}[1]{\bm{#1}}
\newcommand{\tens}[1]{\underline{\underline{#1}}}
\newcommand{\e}[1]{\vect{e}_{#1}}
\newcommand{\dd}{\,\mathrm{d}}
\newcommand{\dv}[2]{\frac{\mathrm{d}#1}{\mathrm{d}#2}}
\newcommand{\Id}{\tens{\mathbb{1}}}
\newcommand{\cm}{\mathrm{cm}}
\newcommand{\ext}{\mathrm{ext}}

% ---------- theorem-like environments ----------
\theoremstyle{plain}
\newtheorem{theorem}{Theorem}[chapter]
\newtheorem{proposition}[theorem]{Proposition}
\newtheorem{corollary}[theorem]{Corollary}
\newtheorem{axiom}[theorem]{Axiom}
\theoremstyle{definition}
\newtheorem{definition}[theorem]{Definition}
\theoremstyle{remark}
\newtheorem{remark}[theorem]{Remark}

% Examples keep the numbering used in the lecture.
\newenvironment{example}[2]{%
  \par\addvspace{\bigskipamount}\noindent
  \textbf{Example\ifx&#1&\else~#1\fi}\ \textit{(#2).}\ \ignorespaces}%
  {\hfill$\blacktriangleleft$\par\addvspace{\bigskipamount}}

% Lecture markers in the margin
\newcommand{\lecture}[2]{\marginpar{\raggedright\footnotesize\textsc{Lecture #1}\\\textit{#2}}}

\title{\Huge\bfseries Mechanics III\\[0.6em]\Large\mdseries Dynamics of Particles, Systems of Particles and Rigid Bodies}
\author{Lecture notes, written up from the handwritten notes of Lectures 1--13}
\date{Autumn Semester 2026}

\begin{document}
\maketitle

\chapter*{Preface}
These notes cover the first thirteen lectures of Mechanics III. They follow the order and the numbering of the
lectures (sections, examples), but whereas the handwritten notes taken in class are necessarily terse, the
goal here is to tell the whole story: every result is motivated, every formula is derived, and every example is
worked through to the end with its physical interpretation.

The course is built around a small number of balance laws. For a single particle, Newton's second law -- the
\emph{balance of linear momentum} -- is the only physical input; the \emph{balance of angular momentum} and
the \emph{work--energy balance} are consequences of it, but each is the most convenient tool in a different
kind of problem. The same three laws are then extended to systems of particles and finally to rigid bodies,
where the angular momentum balance becomes genuinely independent and the moment of inertia tensor appears.
Throughout, the most important practical skill is to choose the right balance law and the right reference
point, so that unknown forces drop out.

\paragraph{Notation.}
Vectors are printed in bold, $\vect r$ (in handwriting and on the blackboard they are underlined,
$\underline r$); second-order tensors are double underlined, $\tens I$. For a vector $\vect r$, the column of its
components in a frame $C$ is written $[\vect r]_C$. A dot denotes the time derivative, $\dot x=\mathrm dx/\mathrm dt$.
We use the abbreviations DOF (degree of freedom), FBD (free-body diagram), LMB, AMB and WEB for the balances of
linear momentum, angular momentum and the work--energy balance, ICR for the instantaneous centre of rotation, IAR
for the instantaneous axis of rotation, and cm for the centre of mass.

\tableofcontents

% =====================================================================
\chapter{Dynamics of a Single Particle}
% =====================================================================
\lecture{1}{16.09.2026}
Dynamics is the study of motion and its causes. We begin with the simplest possible model of a body, the
particle, and develop the two halves of the theory one after the other. \emph{Kinematics}
(\cref{sec:kinematics}) is the geometry of motion: how to describe positions, velocities and accelerations,
without asking why the body moves the way it does. \emph{Kinetics} (\cref{sec:kinetics}) connects the motion to
its cause, the forces, through the balance laws of mechanics.

\begin{definition}
A \emph{particle} is an object of negligible size and shape. Mathematically it is a point carrying a mass $m$.
\end{definition}

Whether this idealisation is appropriate depends on the question being asked, not on the size of the body:
the earth is a particle when we compute its orbit around the sun, but certainly not when we discuss its
rotation. What the definition rules out is precisely that: the orientation of a particle is meaningless, so only
its position has to be described.

% =====================================================================
\section{Kinematics}\label{sec:kinematics}

\subsection{Cartesian reference frames}

To describe a position we need a reference. Throughout the course the basic reference is a Cartesian frame.

\begin{definition}
A \emph{Cartesian reference frame} $C$ consists of a fixed origin $O$ and a fixed orthonormal basis
$\e1,\e2,\e3$, that is, $\e i\cdot\e j=\delta_{ij}$ and $\dot{\vect e}_i=\vect 0$.
\end{definition}

The word ``fixed'' is essential and will become important in \cref{sec:polar}, where we meet bases that move.
The position of the particle at time $t$ is the vector from $O$ to the particle,
\begin{equation}\label{eq:position}
  \vect r(t)=\sum_{i=1}^{3}x_i(t)\,\e i ,
\end{equation}
where the scalar functions $x_i(t)$ are its \emph{coordinates}.

\begin{figure}[ht]
\centering
\begin{tikzpicture}[scale=1.1,>=Stealth]
  \draw[->,thick] (0,0)--(1.4,0) node[right]{$\e1$};
  \draw[->,thick] (0,0)--(0,1.4) node[above]{$\e2$};
  \draw[->,thick] (0,0)--(-0.8,-0.6) node[below left]{$\e3$};
  \fill (0,0) circle (1.2pt) node[below right]{$O$};
  \draw[->,thick] (0,0)--(2.2,1.6) node[midway,above left]{$\vect r(t)$};
  \fill (2.2,1.6) circle (2pt);
  \draw[dashed,gray] (0.8,1.85) .. controls (1.6,1.5) and (2.6,1.9) .. (3.4,1.2) node[right]{trajectory};
\end{tikzpicture}
\caption{A particle in a Cartesian frame.}
\end{figure}

Each coordinate that can be chosen independently is called a \emph{degree of freedom} (DOF). For a particle that
is free to move, all its coordinates are independent; and since the orientation of a particle does not matter,
there are no further DOFs.

\begin{proposition}
A free particle in $d$ dimensions has $d$ degrees of freedom.
\end{proposition}

\paragraph{Velocity and acceleration.}
The velocity is the rate of change of the position, and the acceleration the rate of change of the velocity.
Because the basis vectors of a Cartesian frame are constant, the time derivative acts only on the coordinates:
\begin{equation}
  \vect v=\dv{\vect r}{t}=\sum_{i=1}^{d}\dot x_i\,\e i,\qquad
  \vect a=\dv{\vect v}{t}=\sum_{i=1}^{d}\ddot x_i\,\e i .
\end{equation}
It is worth pausing on the geometric content of these definitions. During an infinitesimal time step $\dd t$
the particle moves by $\dd\vect r=\vect v\dd t$ and its velocity changes by $\dd\vect v=\vect a\dd t$. The first
relation says that the velocity points in the direction of $\dd\vect r$: it is \emph{tangent to the trajectory},
pointing where the particle is going. The second relation says that the acceleration points in the direction of the
\emph{change} of the velocity. In general this has nothing to do with the direction of motion: a car driving
through a bend at constant speed accelerates sideways (\cref{fig:drdv}).

\begin{figure}[ht]
\centering
\begin{tikzpicture}[scale=1.1,>=Stealth]
  \coordinate (O) at (0,0); \coordinate (A) at (-0.6,2.0); \coordinate (B) at (0.6,2.25);
  \draw[dashed,gray] (-1.6,1.4) .. controls (-0.8,2.0) and (0.4,2.3) .. (1.6,2.7);
  \draw[->,thick] (O)--(A) node[midway,left]{$\vect r(t)$};
  \draw[->,thick] (O)--(B) node[midway,right]{$\vect r(t+\dd t)$};
  \draw[->,very thick] (A)--(B) node[midway,above]{$\dd\vect r$};
  \begin{scope}[xshift=4.2cm,yshift=1.2cm]
    \draw[->,thick] (0,0)--(1.5,0) node[below]{$\vect v(t)$};
    \draw[->,thick] (0,0)--(1.6,0.6) node[above]{$\vect v(t+\dd t)$};
    \draw[->,thick,gray] (1.5,0)--(1.6,0.6);
    \node[right] at (1.6,0.15) {$\vect a\dd t$};
  \end{scope}
\end{tikzpicture}
\caption{The velocity is tangent to the trajectory; the acceleration points in the direction of the change of velocity.}
\label{fig:drdv}
\end{figure}

The \emph{speed} is the magnitude of the velocity, $v=|\vect v|=\sqrt{\vect v\cdot\vect v}$. Collecting
coordinates in columns we obtain the component form
\begin{equation}
  [\vect r]_C=\begin{pmatrix}x_1\\x_2\\x_3\end{pmatrix},\qquad
  [\vect v]_C=\begin{pmatrix}\dot x_1\\\dot x_2\\\dot x_3\end{pmatrix},\qquad
  [\vect a]_C=\begin{pmatrix}\ddot x_1\\\ddot x_2\\\ddot x_3\end{pmatrix}.
\end{equation}
The subscript is a reminder that a vector is a geometric object independent of any frame, whereas its components
are not: the same vector has different columns in different frames.

\paragraph{From acceleration to trajectory.}
In kinetics, the balance laws will typically tell us the acceleration; the task is then to reconstruct the motion.
The following examples show how this is done in the most common situations: the acceleration is known as a function
of time, of velocity, or of position.

\begin{example}{1.1}{Finding $\vect r(t)$ from a known $\vect a(t)$}
Starting from $\vect a=\mathrm d\vect v/\mathrm dt$ we separate the differentials and integrate from the initial
instant $t_0$ to the current time $t$:
\[
  \int_{t_0}^{t}\vect a(\tau)\dd\tau=\int_{\vect v(t_0)}^{\vect v(t)}\dd\vect v
  \qquad\Longrightarrow\qquad
  \vect v(t)=\vect v(t_0)+\int_{t_0}^{t}\vect a(\tau)\dd\tau .
\]
Using definite integrals with the initial state as the lower limit has the advantage that no integration constant
has to be determined afterwards. In exactly the same way, $\vect v=\mathrm d\vect r/\mathrm dt$ gives
\[
  \vect r(t)=\vect r(t_0)+\int_{t_0}^{t}\vect v(\tau)\dd\tau .
\]
As an application, consider free fall. With $\e2$ pointing upwards, the acceleration due to gravity is
$\vect a=-g\e2=\text{const}$. Taking $t_0=0$,
\[
  \vect v(t)=\vect v(0)-gt\,\e2,\qquad
  \vect r(t)=\vect r(0)+\vect v(0)\,t-\frac{gt^2}{2}\,\e2 ,
\]
or in components
\[
  \begin{pmatrix}x_1(t)\\x_2(t)\\x_3(t)\end{pmatrix}
  =\begin{pmatrix}x_1(0)\\x_2(0)\\x_3(0)\end{pmatrix}
  +\begin{pmatrix}\dot x_1(0)\\\dot x_2(0)\\\dot x_3(0)\end{pmatrix}t
  +\begin{pmatrix}0\\-gt^2/2\\0\end{pmatrix}.
\]
Since the equation of motion is of second order, the motion is determined only once \emph{two} initial conditions
are given, the initial position $\vect r(0)$ and the initial velocity $\vect v(0)$.
\end{example}

\begin{example}{}{Clicker question: which particle lands first?}
Two identical particles are launched at the same time. Particle 1 starts from the ground vertically upwards with
speed $v_0$; particle 2 is released from rest at height $h$. Both move with the constant acceleration $g$
downwards. From what height must particle 2 be released for particle 1 to land first?

By the previous example, the height of particle 1 is
$x_2^{(1)}(t)=v_0t-\tfrac12gt^2=t\big(v_0-\tfrac12gt\big)$. It vanishes at $t=0$ (the launch) and at
$t^{(1)}=2v_0/g$ (the landing). Particle 2 has $x_2^{(2)}(t)=h-\frac12gt^2$, which vanishes at
$t^{(2)}=\sqrt{2h/g}$. Particle 1 lands first if $t^{(1)}<t^{(2)}$; both sides are positive, so we may square:
\[
  \frac{4v_0^2}{g^2}<\frac{2h}{g}\quad\Longleftrightarrow\quad h>\frac{2v_0^2}{g}.
\]
As a plausibility check, particle 1 rises to $v_0^2/(2g)$, so particle 2 has to start four times as high.
\end{example}

\begin{example}{1.5}{Finding the motion from $a(v)$ or $a(x)$}
\emph{(1) Drag.} A particle moving through a medium is slowed down by it. For slow motion the drag is well
described by Stokes' law, which makes the acceleration proportional to the velocity, $\vect a=-k\vect v$ with
$k>0$. In one dimension ($\vect r=x\e1$, $\vect v=v\e1$, $\vect a=a\e1$) we have $a=a(v)=-kv$. The time does not
appear explicitly, so we separate variables in $a(v)=\mathrm dv/\mathrm dt$:
\[
  \int_{t_0}^{t}\dd\tau=\int_{v(t_0)}^{v(t)}\frac{\dd v}{a(v)}=-\frac1k\int_{v(t_0)}^{v(t)}\frac{\dd v}{v}
  \qquad\Longrightarrow\qquad t-t_0=-\frac1k\ln\frac{v(t)}{v(t_0)} ,
\]
that is, $v(t)=v(t_0)\,e^{-k(t-t_0)}$. The velocity decays exponentially; one more integration gives
$x(t)=x(t_0)+\frac{v(t_0)}{k}\big(1-e^{-k(t-t_0)}\big)$, so the particle comes to rest after the finite distance
$v(t_0)/k$, although it needs infinitely long to do so.

\emph{(2) Rocket launch.} Far from the surface of the earth, the gravitational acceleration depends on the altitude:
$a=a(x)$. Now the trick is to eliminate the time by the chain rule,
\[
  a(x)=\dv{v}{t}=\dv{v}{x}\,\dv{x}{t}=v\,\dv{v}{x},
\]
which separates into $v\dd v=a(x)\dd x$ and integrates to
\[
  \tfrac12\big(v(x)^2-v(x_0)^2\big)=\int_{x_0}^{x}a(\xi)\dd\xi .
\]
This gives the velocity as a function of position, $v=v(x)$. To obtain the position as a function of time, we
separate once more: $\dd t=\dd x/v(x)$, hence $t-t_0=\int_{x_0}^{x}\dd\xi/v(\xi)$, which can be inverted to give
$x(t)$. The relation $\frac12\mathrm d(v^2)=a\dd x$ obtained here is a first glimpse of the work--energy balance of
\cref{sec:web}.
\end{example}

\subsection{Constrained motion}
Very often a particle is not free: a bead on a wire, a block on a table, a mass on a string. Each such condition is a
\emph{kinematic constraint}, a relation that the coordinates must satisfy at all times. Each independent constraint
allows us to express one coordinate in terms of the others and therefore removes one degree of freedom.

\begin{proposition}
A particle in $d$ dimensions subject to $k$ independent kinematic constraints has $d-k$ degrees of freedom.
\end{proposition}

\begin{example}{1.6}{(Un-)constrained motion}
A particle moving freely in space, $\vect r=\sum_{i=1}^dx_i\e i$, has $d$ DOFs. A particle on a rigid plane with
normal $\vect n$ must satisfy $\vect r\cdot\vect n=\text{const}$; in three dimensions it therefore has $3-1=2$ DOFs,
and choosing $\e1,\e2$ in the plane we can write $\vect r=x_1\e1+x_2\e2$. A particle attached to a fixed point by a
rigid link of length $L$ satisfies $|\vect r|=L$ and has $d-1$ DOFs; it moves on a circle in two dimensions and on a
sphere in three. Finally, a particle confined to a tube or a track has a single DOF: its position is completely
described by the path length $s$ travelled along the track.
\end{example}

Constraints do not enforce themselves. Whatever keeps the particle on the plane, on the circle or in the tube is a
force, and so we note the following principle, which will be used in every problem of kinetics:
\emph{every constraint gives rise to a reaction force} -- the normal force of a surface, the tension in a string --
which is exactly the force needed to enforce the constraint. Conversely, a constraint that changes the direction of
the velocity produces an acceleration (even at constant speed), and by Newton's law this requires a force.

% ---------------------------------------------------------------------
\subsection{Kinematics in polar coordinates}\label{sec:polar}
\lecture{2}{21.09.2026}
Many motions are naturally described by a distance and an angle: a mass on a string, a planet around the sun, a
bead on a rotating rod. In two dimensions we therefore describe the position by the distance $r(t)$ from the origin
and the angle $\varphi(t)$ measured from $\e1$:
\begin{equation}
  \vect r(t)=r\cos\varphi\,\e1+r\sin\varphi\,\e2 .
\end{equation}
We could differentiate this directly. The product and chain rules give
\begin{equation*}
  \vect v=\big(\dot r\cos\varphi-r\dot\varphi\sin\varphi\big)\e1+\big(\dot r\sin\varphi+r\dot\varphi\cos\varphi\big)\e2,
\end{equation*}
and differentiating once more produces a rather unpleasant expression for the acceleration, whose physical meaning
is completely hidden. The remedy is to use a basis adapted to the motion.

\begin{definition}
The \emph{polar basis} associated with the angle $\varphi$ is
\begin{equation}
  \e r(t)=\cos\varphi\,\e1+\sin\varphi\,\e2,\qquad \e\varphi(t)=-\sin\varphi\,\e1+\cos\varphi\,\e2 .
\end{equation}
\end{definition}

The vector $\e r$ points radially outwards towards the particle, and $\e\varphi$ is obtained from it by a quarter
turn in the direction of increasing $\varphi$. One checks immediately that $|\e r|=|\e\varphi|=1$ and
$\e r\cdot\e\varphi=0$, so the polar basis is orthonormal. The crucial difference to a Cartesian basis is that it
co-rotates with the particle and is therefore \emph{time dependent}.

\begin{figure}[ht]
\centering
\begin{tikzpicture}[scale=1.3,>=Stealth]
  \draw[->,thick] (0,0)--(1.2,0) node[right]{$\e1$};
  \draw[->,thick] (0,0)--(0,1.2) node[above]{$\e2$};
  \draw[thick] (0,0)--(2.4,1.6);
  \fill (2.4,1.6) circle (2pt);
  \node at (1.0,1.0) {$r$};
  \draw[->,thick] (2.4,1.6)--++(0.83,0.555) node[right]{$\e r$};
  \draw[->,thick] (2.4,1.6)--++(-0.555,0.83) node[above]{$\e\varphi$};
  \draw (0.7,0) arc (0:33.7:0.7); \node at (0.9,0.25){$\varphi$};
\end{tikzpicture}
\caption{The co-rotating polar basis.}
\end{figure}

Since the basis moves, we need its time derivatives. Differentiating with the chain rule,
\begin{align}
  \dot{\vect e}_r&=-\sin\varphi\,\dot\varphi\,\e1+\cos\varphi\,\dot\varphi\,\e2=\dot\varphi\,\e\varphi,\label{eq:der-er}\\
  \dot{\vect e}_\varphi&=-\cos\varphi\,\dot\varphi\,\e1-\sin\varphi\,\dot\varphi\,\e2=-\dot\varphi\,\e r.\label{eq:der-ephi}
\end{align}
These two identities are all we need. They are easy to remember geometrically: a unit vector cannot change its
length, so its derivative is perpendicular to it, and its magnitude is the rate $\dot\varphi$ at which the vector
turns.

With the polar basis the position takes the compact form $\vect r=r\,\e r$. Differentiating with the product rule
and using \eqref{eq:der-er},
\begin{equation}\label{eq:polar-v}
  \vect v=\dot r\,\e r+r\,\dot{\vect e}_r=\dot r\,\e r+r\dot\varphi\,\e\varphi .
\end{equation}
The two terms have a clear meaning: the radial component $\dot r$ comes from the change of distance, the
circumferential component $r\dot\varphi$ from the rotation. We call $\omega\coloneqq\dot\varphi$ the \emph{angular
velocity} (unit $\si{\radian\per\second}$) and $\dot\omega=\ddot\varphi$ the \emph{angular acceleration}
($\si{\radian\per\second\squared}$). Differentiating \eqref{eq:polar-v} once more, now using both
\eqref{eq:der-er} and \eqref{eq:der-ephi},
\begin{align*}
  \vect a&=\ddot r\,\e r+\dot r\,\dot{\vect e}_r+\dot r\dot\varphi\,\e\varphi+r\ddot\varphi\,\e\varphi+r\dot\varphi\,\dot{\vect e}_\varphi\\
  &=\ddot r\,\e r+\dot r\dot\varphi\,\e\varphi+\dot r\dot\varphi\,\e\varphi+r\ddot\varphi\,\e\varphi-r\dot\varphi^2\,\e r,
\end{align*}
and collecting terms we arrive at
\begin{equation}\label{eq:polar-a}
  \vect a=\big(\ddot r-r\dot\varphi^2\big)\e r+\big(2\dot r\dot\varphi+r\ddot\varphi\big)\e\varphi .
\end{equation}
Each term has a name and a meaning. The term $-r\dot\varphi^2$ is the \emph{centripetal acceleration}; it always
points towards the origin and is present whenever the particle rotates, even at constant $r$ and $\dot\varphi$.
The term $2\dot r\dot\varphi$ is the \emph{Coriolis acceleration}: if you walk outwards on a rotating platform, your
circumferential speed $r\dot\varphi$ grows with $r$, so you must be accelerated sideways. Finally, $r\ddot\varphi$
is the tangential acceleration caused by a change of the rotation rate.

\begin{remark}
Note that $v_r=\dot r$, but $a_r\neq\dot v_r$, and likewise for the $\varphi$-components. Because the basis is time
dependent, the components of the acceleration are \emph{not} the derivatives of the components of the velocity.
Forgetting this is one of the most common mistakes in polar coordinates.
\end{remark}

\begin{example}{1.7}{Particle on a taut, inextensible string}
The particle moves on a circle of radius $L$, so $r=L=\text{const}$ and $\dot r=\ddot r=0$. Equations
\eqref{eq:polar-v} and \eqref{eq:polar-a} reduce to
\[
  \vect v=L\dot\varphi\,\e\varphi,\qquad\vect a=-L\dot\varphi^2\,\e r+L\ddot\varphi\,\e\varphi ,
\]
and the speed is $|\vect v|=L|\omega|$. For uniform rotation, $\omega=\text{const}$, only the centripetal term
survives: $\vect a=-L\omega^2\e r=-(v^2/L)\e r$. The particle accelerates although its speed is constant, because
the direction of its velocity keeps changing.
\end{example}

\begin{example}{}{Clicker question: moving radially on a rotating disk}
A disk rotates with constant angular velocity $\omega$, and a particle moves radially inwards on it with constant
speed $v_0$; thus $\dot r=-v_0$, $\ddot r=0$, $\dot\varphi=\omega$, $\ddot\varphi=0$. Then
\[
  \vect v=-v_0\e r+r\omega\,\e\varphi,\qquad\vect a=-r\omega^2\e r-2v_0\omega\,\e\varphi,\qquad
  |\vect a|=\sqrt{(r\omega^2)^2+(2v_0\omega)^2}.
\]
Since $v_0$ and $\omega$ are constant, the magnitude of the acceleration is smallest at the centre $r=0$, where it
equals $2|v_0\omega|$. The Coriolis part never vanishes.
\end{example}

\subsection{Kinematics of space curves}
When the path of the particle is prescribed -- a road, a rail, a roller coaster track -- the most natural coordinate
is the \emph{path length} $s(t)$ measured along the curve from some reference point, and the most natural basis is
attached to the curve itself. The \emph{unit tangent} $\e t=\vect v/|\vect v|$ points in the direction of motion,
so that
\begin{equation}
  \vect v=v\,\e t=\dot s\,\e t ,
\end{equation}
where we have used that the speed is the rate at which path length is covered, $v=|\vect v|=\dot s$.

The acceleration follows from the product rule, $\vect a=\ddot s\,\e t+\dot s\,\dot{\vect e}_t$, and the question is
how fast the tangent turns. Locally, every smooth curve can be approximated by a circle, the \emph{osculating
circle}, whose radius $\rho$ is the \emph{radius of curvature}. Along a circle of radius $\rho$ a path element
$\dd s$ corresponds to a change of direction $\dd\theta=\dd s/\rho$, so the tangent changes at the rate
$\mathrm d\e t/\mathrm ds=\e n/\rho$, where the \emph{unit normal} $\e n$ points towards the centre of the
osculating circle. Hence $\dot{\vect e}_t=(\dot s/\rho)\e n$ and
\begin{equation}\label{eq:path-a}
  \vect a=\ddot s\,\e t+\frac{v^2}{\rho}\,\e n .
\end{equation}
The first term changes the speed along the path; the second, always directed towards the centre of curvature,
changes the direction. On a straight piece of track $\rho\to\infty$, the normal term disappears and
$\vect a=\ddot s\,\e t$.

\begin{example}{1.10}{Circular motion}
On a circle of radius $R$ we have $\rho=R$ and $s=R\varphi$, so \eqref{eq:path-a} gives
$\vect v=R\dot\varphi\,\e t$ and $\vect a=R\ddot\varphi\,\e t+(v^2/R)\,\e n$. Polar coordinates with $r=R$ gave
$\vect a=R\ddot\varphi\,\e\varphi-(v^2/R)\,\e r$. The two results agree: $\e t=\e\varphi$, and since $\e n$ points
towards the centre while $\e r$ points away from it, $\e n=-\e r$.
\end{example}

Formula \eqref{eq:path-a} has an immediate practical consequence for track design. If a straight track is joined
directly to a circular arc of radius $R$, the curvature jumps from $0$ to $1/R$ at the junction, and so does the
normal acceleration $v^2/\rho$ -- passengers are hit by a sudden sideways jolt. Good designs insert a transition
curve along which the radius of curvature decreases gradually (for instance a clothoid, whose curvature grows
linearly with $s$), so that $v^2/\rho$ builds up smoothly; in a looping one aims to keep $v^2/\rho$ roughly
constant, which is why real loopings are not circular.

% =====================================================================
\section{Kinetics}\label{sec:kinetics}
\lecture{3}{23.09.2026}
We now turn to the causes of motion. The central quantity is the momentum of the particle.

\begin{definition}
The \emph{linear momentum} of a particle of mass $m$ and velocity $\vect v$ is $\vect p=m\vect v$.
\end{definition}

\subsection{Balance of linear momentum}
The fundamental law of mechanics, Newton's second law, states how the momentum changes.

\begin{axiom}[Balance of linear momentum]
In an inertial frame, the rate of change of the momentum of a particle equals the net force acting on it:
\begin{equation}\label{eq:lmb}
  \sum_i\vect F_i=\dv{\vect p}{t}=\dv{}{t}(m\vect v).
\end{equation}
\end{axiom}

Several familiar statements are special cases. In \emph{statics} nothing moves, $\vect v=\vect 0$, so $\vect p=\vect 0$
and the forces must balance, $\sum_i\vect F_i=\vect 0$. Conversely, if the net force vanishes, the momentum is
constant and so is the velocity of a particle of constant mass: it moves uniformly along a straight line. This is
\emph{Newton's first axiom}. If the mass is constant, $\dot m=0$, then $\dot{\vect p}=m\dot{\vect v}$ and we
obtain \emph{Newton's second axiom} in its usual form,
\begin{equation}
  \sum_i\vect F_i=m\vect a .
\end{equation}
It is the acceleration, not the velocity, that is determined by the forces. Finally, for a massless object ($m=0$)
the momentum is always zero, so the forces on it must balance \emph{even when it accelerates}. This is why we may
treat light ropes, pulleys and links as if they were in static equilibrium.

The third axiom completes the picture.

\begin{axiom}[Newton's third axiom, ``actio = reactio'']
If a body $A$ exerts a force $\vect F$ on a body $B$, then $B$ exerts the force $-\vect F$ on $A$.
\end{axiom}

This axiom is what makes free-body diagrams consistent. To apply \eqref{eq:lmb} to one body of a mechanism, we
cut it free from its surroundings and replace every contact by the corresponding force; the same force then appears
with the opposite sign on the body on the other side of the cut.

\begin{example}{}{Free-body diagram of a mass on a block}
A small mass $m$ lies on a block of mass $M$, which rests on the ground; $\e2$ points up. Cutting the mass free, it
experiences its weight $-mg\e2$ and the contact force $N_1\e2$ from the block. The block experiences its weight
$-Mg\e2$, the reaction $-N_1\e2$ of the mass, and the force $N_2\e2$ from the ground. (The ground actually acts
through a distributed pressure, but only its resultant enters the momentum balance.) Applying \eqref{eq:lmb} in the
vertical direction to each body,
\[
  N_1-mg=ma_1,\qquad N_2-N_1-Mg=Ma_2 .
\]
\end{example}

\begin{example}{1.12}{Block sliding on frictionless ground}
The block can only move horizontally, so $\vect r=x_1\e1$, $\vect v=\dot x_1\e1$, $\vect a=\ddot x_1\e1$. The friction
coefficient is $\mu=0$, so the only forces are gravity and the normal force, and the momentum balance reads
\[
  (N-mg)\,\e2=m\ddot x_1\,\e1 .
\]
The left-hand side has only an $\e2$-component and the right-hand side only an $\e1$-component, so both must
vanish: $N=mg$ and $\ddot x_1=0$. The block slides on with constant velocity.
\end{example}

Real surfaces resist sliding. In the Coulomb model of friction, as long as the surfaces stick, the friction force
$\vect R$ adjusts itself to whatever value is needed to prevent sliding, but only up to a limit,
$|\vect R|\le\mu_0|\vect N|$, where $\mu_0$ is the static friction coefficient. Once the surfaces slide, the
friction force has a fixed magnitude and opposes the sliding velocity:
\begin{equation}
  \vect R=-\mu|\vect N|\,\frac{\vect v}{|\vect v|},
\end{equation}
with the dynamic friction coefficient $\mu$.

\begin{example}{1.13}{Sliding down a rough incline}
\begin{center}
\begin{tikzpicture}[scale=1.0,>=Stealth]
  \draw[thick] (0,0)--(5,0)--(0,2.9)--cycle;
  \fill[pattern=north east lines] (0,0) rectangle (5,-0.15);
  \begin{scope}[shift={(2.0,1.74)},rotate=-30]
    \draw[fill=gray!20,thick] (-0.5,0) rectangle (0.5,0.6);
    \draw[->,thick] (0,0.6)--(0,1.4) node[above]{$N$};
    \draw[->,thick] (-0.5,0.3)--(-1.3,0.3) node[left]{$R$};
    \draw[->] (1.2,0.3)--(2.0,0.3) node[right]{$\e1$};
    \draw[->] (1.2,0.3)--(1.2,1.1) node[above]{$\e2$};
  \end{scope}
  \draw[->,thick] (2.15,2.0)--(2.15,0.9) node[below]{$mg$};
  \draw (4.2,0) arc (180:150:0.8); \node at (3.9,0.2) {$\alpha$};
\end{tikzpicture}
\end{center}
We choose $\e1$ down the slope and $\e2$ normal to it, so that the kinematics is again one-dimensional,
$\vect a=\ddot x_1\e1$. The block slides downwards, so the dynamic friction force points up the slope,
$\vect R=-R\e1$ with $R=\mu N$. Decomposing the weight into the components $mg\sin\alpha$ along $\e1$ and
$-mg\cos\alpha$ along $\e2$, the momentum balance is
\[
  (mg\sin\alpha-R)\,\e1+(N-mg\cos\alpha)\,\e2=m\ddot x_1\,\e1 .
\]
The $\e2$-component gives the normal force, $N=mg\cos\alpha$, and the $\e1$-component together with $R=\mu N$
gives the acceleration
\[
  \ddot x_1=g\big(\sin\alpha-\mu\cos\alpha\big).
\]
If $\tan\alpha<\mu$ this is negative: the block decelerates and eventually comes to rest.
\end{example}

\begin{example}{1.14}{Projectile motion}
If gravity is the only force, the momentum balance $m\vect a=m\vect g$ gives $\vect a=\vect g$, independently of
the mass. The horizontal velocity therefore stays exactly the same during the whole flight, while the vertical
motion is the free fall of Example 1.1; the trajectory is a parabola.
\end{example}

\begin{example}{1.15}{Particle on a string}
A mass $m$ on a string of constant length $L$ rotates about a fixed point with constant angular velocity $\omega$.
From Example 1.7, $\vect a=-L\omega^2\e r$. The only force (we ignore gravity) is the string force, which pulls
towards the centre, $-S\e r$. The momentum balance $-S\e r=-mL\omega^2\e r$ thus gives $S=mL\omega^2$: the string
has to supply exactly the centripetal force.
\end{example}

% ---------------------------------------------------------------------
\subsection{Balance of angular momentum}
For rotational motion it is useful to consider the moment of the momentum about a point.

\begin{definition}
The \emph{angular momentum} of a particle $P$ with respect to a reference point $B$ is
\begin{equation}
  \vect H_B=\vect r_{BP}\times\vect p,\qquad\vect r_{BP}=\vect r_P-\vect r_B,\quad\vect p=m\vect v_P .
\end{equation}
\end{definition}

\begin{theorem}[Balance of angular momentum]
Let $\vect M_B=\vect r_{BP}\times\sum_i\vect F_i$ be the net torque about $B$. Then
\begin{equation}\label{eq:amb-particle}
  \vect M_B=\dot{\vect H}_B+\vect v_B\times\vect p .
\end{equation}
\end{theorem}
\begin{proof}[Derivation]
By the product rule, $\dot{\vect H}_B=\dot{\vect r}_{BP}\times\vect p+\vect r_{BP}\times\dot{\vect p}$. In the
first term $\dot{\vect r}_{BP}=\vect v_P-\vect v_B$, and $\vect v_P\times\vect p=m\,\vect v_P\times\vect v_P=\vect 0$,
so it equals $-\vect v_B\times\vect p$. In the second term we use the momentum balance $\dot{\vect p}=\sum_i\vect F_i$,
which turns it into the torque $\vect M_B$. Hence $\dot{\vect H}_B=-\vect v_B\times\vect p+\vect M_B$.
\end{proof}

For a single particle the angular momentum balance contains no new physics -- we have just derived it from the
momentum balance. Its value lies in its convenience: forces whose line of action passes through $B$ have no torque
about $B$ and simply drop out, so a clever choice of $B$ eliminates unknown reaction forces.

The extra term $\vect v_B\times\vect p$ disappears if the reference point is fixed, $\vect v_B=\vect 0$, or if it moves
parallel to the particle. In that case the balance takes the simple form $\vect M_B=\dot{\vect H}_B$, and the
reference point should be chosen in this way whenever possible. In particular, if there is no torque about such a
point, the angular momentum about it is conserved. In statics, and for a massless particle, both $\vect H_B$ and
$\vect p$ vanish and the balance reduces to $\vect M_B=\vect 0$.

\paragraph{Plane motion about a fixed point.}
In two dimensions, with $B$ fixed at the origin of polar coordinates, $\vect r_{BP}=r\e r$ and
$\vect v_P=v_r\e r+v_\varphi\e\varphi$. Since $\e r\times\e r=\vect 0$ and $\e r\times\e\varphi=\e3$,
\[
  \vect H_B=m\,r\e r\times\big(v_r\e r+v_\varphi\e\varphi\big)=mrv_\varphi\,\e3=mr^2\omega\,\e3 .
\]
The angular momentum always points out of the plane, and writing $\vect M_B=M_B\e3$ the balance becomes a scalar
equation,
\begin{equation}
  M_B=\dv{}{t}\big(mr^2\omega\big).
\end{equation}
If, moreover, the distance is constant, $r=R$, this reduces to
\begin{equation}
  M_B=mR^2\dot\omega=I_B\ddot\varphi,\qquad I_B=mR^2 .
\end{equation}
The quantity $I_B$ is the \emph{mass moment of inertia} about $B$. It plays the same role for rotation that the
mass plays for translation -- compare $F=ma$ with $M_B=I_B\ddot\varphi$ -- and measures the resistance of the body
against angular acceleration. We shall meet it again, in much greater generality, in \cref{ch:rigid}.

% ---------------------------------------------------------------------
\subsection{Work--energy balance}\label{sec:web}
\lecture{4}{28.09.2026}
The momentum balance is a differential equation in time. In many problems, however, we are not interested in
\emph{when} something happens but only in the relation between the state of motion at two different positions --
how fast is the particle at the bottom of the hill? For such questions an integrated form of the momentum balance
is far more efficient.

\begin{theorem}[Work--energy balance]
For a particle of constant mass moving from $\vect r_1=\vect r(t_1)$ to $\vect r_2=\vect r(t_2)$,
\begin{equation}\label{eq:web}
  T(t_2)-T(t_1)=W_{12},\qquad T=\frac m2|\vect v|^2,\qquad W_{12}=\int_{\vect r_1}^{\vect r_2}\sum_i\vect F_i\cdot\dd\vect r .
\end{equation}
$T$ is the \emph{kinetic energy} and $W_{12}$ the \emph{work} done by the forces on the particle: the change of
kinetic energy equals the work done.
\end{theorem}
\begin{proof}[Derivation]
Take the scalar product of the momentum balance with $\dd\vect r$ and integrate along the path:
\[
  \int_{\vect r_1}^{\vect r_2}\dv{}{t}(m\vect v)\cdot\dd\vect r=\int_{\vect r_1}^{\vect r_2}\sum_i\vect F_i\cdot\dd\vect r=W_{12}.
\]
On the left we substitute $\dd\vect r=\vect v\dd t$ and change the integration variable:
\[
  m\int_{t_1}^{t_2}\dv{\vect v}{t}\cdot\vect v\dd t=m\int_{\vect v(t_1)}^{\vect v(t_2)}\vect v\cdot\dd\vect v
  =\frac m2|\vect v(t_2)|^2-\frac m2|\vect v(t_1)|^2 .\qedhere
\]
\end{proof}

An important observation is that a force perpendicular to the motion, $\vect F\perp\dd\vect r$, does no work. The
normal force of a fixed surface and the string force of a pendulum are of this kind, and so they disappear from the
work--energy balance altogether. This is what makes it so effective.

\begin{example}{1.21}{Frictional sliding}
A block with initial velocity $v_0$ slides on rough ground. Where does it come to rest? The kinetic energy is
$\frac12mv_0^2$ initially and zero at the end. Of the forces acting on the block, gravity and the normal force
are perpendicular to the motion $\dd\vect r=\dd x_1\e1$ and do no work; only friction, $-\mu N\e1=-\mu mg\e1$, does.
Thus $W_{12}=-\mu mg\,\Delta x$, and \eqref{eq:web} gives
\[
  0-\tfrac12mv_0^2=-\mu mg\,\Delta x\qquad\Longrightarrow\qquad\Delta x=\frac{v_0^2}{2\mu g}.
\]
We never needed to know how long the block slides.
\end{example}

\paragraph{Conservative forces.}
Computing the work requires, in general, knowledge of the path. For an important class of forces it does not.

\begin{definition}
A force $\vect F$ is \emph{conservative} if there is a function $V(\vect r)$, the \emph{potential energy}, such that
$\vect F=-\mathrm dV/\mathrm d\vect r=-\nabla V$.
\end{definition}

For instance, gravity is conservative: with $\e2$ up, $V(\vect r)=mg\,\vect r\cdot\e2=mgx_2$ gives
$\vect F=-\nabla V=-mg\e2$.

\begin{proposition}
The work of a conservative force depends only on the end points of the path: $W_{12}=V(\vect r_1)-V(\vect r_2)$.
\end{proposition}
\begin{proof}
$W_{12}=\int_{\vect r_1}^{\vect r_2}\vect F\cdot\dd\vect r=-\int_{\vect r_1}^{\vect r_2}\nabla V\cdot\dd\vect r=-\int_{V(\vect r_1)}^{V(\vect r_2)}\dd V$.
\end{proof}

We can check this on a particle sliding down a frictionless curved track from height $h_1$ to $h_2$. The normal
force is perpendicular to the velocity, so only gravity does work,
$W_{12}=\int(-mg\e2)\cdot\dd\vect r=-mg\int\dd x_2=-mg(h_2-h_1)=V_1-V_2$, independently of the shape of the track.

If we now split the forces into conservative forces, which we account for by their potential, and the remaining
non-conservative ones, the work is $W_{12}=V(\vect r_1)-V(\vect r_2)+W_{12}^{\text{nc}}$, and the work--energy
balance takes the form
\begin{equation}\label{eq:web-pot}
  T(t_2)+V(\vect r_2)-\big[T(t_1)+V(\vect r_1)\big]=W_{12}^{\text{nc}},\qquad
  W_{12}^{\text{nc}}=\int_{\vect r_1}^{\vect r_2}\sum_i\vect F_i^{\text{nc}}\cdot\dd\vect r .
\end{equation}
A system in which work is done only by conservative forces is called \emph{conservative}. For it $W^{\text{nc}}_{12}=0$,
and the total mechanical energy is conserved:
\begin{equation}
  T+V=\text{const}.
\end{equation}

\begin{example}{1.22}{Maximum flight height}
A projectile is launched with speed $v_0$ at an angle $\alpha$ to the horizontal; there is no air drag. At the
highest point the vertical velocity vanishes, while the horizontal velocity $v_0\cos\alpha$ is unchanged
(Example 1.14). Conservation of energy between launch and top gives
\[
  \frac m2v_0^2=\frac m2(v_0\cos\alpha)^2+mgx_{2,\max}\qquad\Longrightarrow\qquad
  x_{2,\max}=\frac{v_0^2}{2g}\big(1-\cos^2\alpha\big)=\frac{v_0^2}{2g}\sin^2\alpha .
\]
\end{example}

\begin{example}{1.23}{Car in a looping}
A car, modelled as a particle of mass $m$, starts from rest at height $h_0$ and rolls without friction and without
engine power down a ramp into a circular looping of radius $R$. We ask for the minimal height $h_0$ for which the
car completes the looping, and for the largest force it experiences.

Energy conservation gives the speed at height $h$: $mgh_0=mgh+\frac12mv^2$, so $v^2=2g(h_0-h)$. Requiring $v^2\ge0$
at the top of the looping yields $h_0\ge2R$. This condition is necessary, but not sufficient: with $h_0$ only
slightly above $2R$, the car would still have speed left at the top, but long before reaching it it would lose
contact with the track and fall. The real condition is that the track must keep pushing the car,
$N(\varphi)\ge0$.

To compute $N$ we measure $\varphi$ from the bottom (the entrance) of the looping and use polar coordinates about its
centre, with $\e r$ pointing outwards. The normal force points inwards, the radial component of gravity is
$mg\cos\varphi$, and the radial acceleration is $-R\dot\varphi^2=-v^2/R$. The radial momentum balance therefore reads
\[
  -N(\varphi)+mg\cos\varphi=-m\frac{v^2}{R}\qquad\Longrightarrow\qquad N(\varphi)=mg\cos\varphi+\frac{mv^2}{R}.
\]
Inserting $v^2=2g(h_0-h)$ with $h=R(1-\cos\varphi)$,
\[
  N(\varphi)=mg\left(3\cos\varphi-2+\frac{2h_0}{R}\right).
\]
The normal force is smallest at the top, $\varphi=\pi$, where $N(\pi)=mg(2h_0/R-5)$, and the condition $N(\pi)\ge0$
gives
\[
  h_0\ge\frac52R .
\]
The force is largest at the entrance $\varphi=0$; for the minimal height $h_0=5R/2$ it is $N(0)=mg(3-2+5)=6mg$. The
passengers would feel six times their own weight -- another reason why loopings are not built circular.
\end{example}

\lecture{5}{30.09.2026}
The pendulum provides a good opportunity to compare the three balance laws on one and the same problem.

\begin{example}{1.17}{Pendulum}
A mass $m$ hangs on a massless string of length $R$ from a fixed pivot $B$; $\varphi$ denotes the angle from the
downward vertical. We seek the angular acceleration $\ddot\varphi$. The forces are the string force $-F\e r$ and the
weight, which in the polar basis about $B$ has the components $mg\cos\varphi\,\e r-mg\sin\varphi\,\e\varphi$.

\emph{Momentum balance.} With $r=R$ the acceleration is $-R\dot\varphi^2\e r+R\ddot\varphi\,\e\varphi$, so
\[
  (-F+mg\cos\varphi)\,\e r-mg\sin\varphi\,\e\varphi=m\big(-R\dot\varphi^2\e r+R\ddot\varphi\,\e\varphi\big).
\]
The $\e\varphi$-component gives the equation of motion $\ddot\varphi=-(g/R)\sin\varphi$, and the $\e r$-component the
string force $F=mg\cos\varphi+mR\dot\varphi^2$; note the same structure as the normal force in the looping.

\emph{Angular momentum balance.} About the fixed pivot the string force has no torque, and gravity has the lever
arm $R\sin\varphi$. Thus $-mgR\sin\varphi=mR^2\ddot\varphi$, which is the same equation of motion, obtained without
ever introducing the unknown string force.

\emph{Work--energy balance.} The system is conservative (the string force is perpendicular to the motion). With the
potential measured from the height of the pivot and the mass released from rest at height $-h_0$,
\[
  -mgh_0=-mgR\cos\varphi+\frac m2R^2\dot\varphi^2 .
\]
The left-hand side is constant, so differentiating with respect to time gives
$0=mgR\sin\varphi\,\dot\varphi+mR^2\dot\varphi\ddot\varphi$, i.e.\ once more $\ddot\varphi=-(g/R)\sin\varphi$.
For small angles, $\sin\varphi\approx\varphi$, this is the equation of a harmonic oscillator with angular frequency
$\sqrt{g/R}$.
\end{example}

% ---------------------------------------------------------------------
\subsection{Particle impact}
We finally consider a particle hitting a rigid and smooth wall. ``Rigid'' means that the wall neither moves nor
deforms; ``smooth'' means that it is frictionless and can only push along its normal $\e n$. During the impact the
contact force $\vect F(t)$ is very large but acts only for a very short time $[t_-,t_+]$, so that it is not the
force itself that matters but its time integral.

\begin{definition}
The \emph{impulse} of a force over the interval $[t_-,t_+]$ is $\hat{\vect F}=\int_{t_-}^{t_+}\vect F(t)\dd t$
(unit $\si{\newton\second}$).
\end{definition}

Integrating the momentum balance over the impact gives $m\vect v(t_+)-m\vect v(t_-)=\hat{\vect F}$: the impulse is
the total momentum transferred to the particle. It is helpful to split the impact into two phases at the instant
$t_0$ of maximum compression. During the \emph{compression} phase $[t_-,t_0]$ the particle is slowed down and
deformed; during the \emph{restitution} phase $[t_0,t_+]$ it (partially) regains its original shape and is pushed
away. Denote the corresponding impulses by $\hat F_{\text{comp}}\e n$ and $\hat F_{\text{rest}}\e n$ -- both are along
$\e n$, since the wall is smooth.

Decomposing the velocity into normal and tangential parts, $\vect v=v_n\e n+v_t\e t$, the momentum balance for the
two phases reads
\[
  m\vect v(t_0)-m\vect v(t_-)=\hat F_{\text{comp}}\e n,\qquad m\vect v(t_+)-m\vect v(t_0)=\hat F_{\text{rest}}\e n .
\]
The tangential components receive no impulse and are therefore unchanged,
\begin{equation}
  v_t(t_-)=v_t(t_0)=v_t(t_+).
\end{equation}
At the instant of maximum compression the normal velocity vanishes, $v_n(t_0)=0$, so the normal components give
$-mv_n(t_-)=\hat F_{\text{comp}}$ and $mv_n(t_+)=\hat F_{\text{rest}}$. These are two equations for three unknowns.
The missing information is a property of the material: how much of the compression is recovered.

\begin{definition}
The \emph{coefficient of restitution} is $e=\hat F_{\text{rest}}/\hat F_{\text{comp}}\in[0,1]$. The impact is called
\emph{elastic} if $e=1$ (full recovery), \emph{plastic} if $e=0$ (no recovery), and partially elastic otherwise.
\end{definition}

Dividing the two normal equations by each other, $-v_n(t_+)/v_n(t_-)=e$, that is,
\begin{equation}
  v_n(t_+)=-e\,v_n(t_-).
\end{equation}
In an elastic impact the normal velocity is simply reversed; in a plastic impact the particle stays at the wall and
slides along it with its unchanged tangential velocity; in between, the normal speed is reduced by the factor $e$.
That $e$ cannot exceed one follows from energy considerations: the kinetic energies before and after are
$\frac m2(v_n^2+v_t^2)$ and $\frac m2(e^2v_n^2+v_t^2)$, and for $e>1$ the passive wall would create energy out of
nothing.

\begin{example}{1.26}{Measuring the coefficient of restitution}
A particle is dropped from rest from height $h_-$ onto the ground and rebounds to height $h_+$. The work--energy
balance does not hold \emph{across} the impact, where energy is lost, but it does hold during the free flights
before and after. Hence $mgh_-=\frac m2v_n^2(t_-)$ and $mgh_+=\frac m2v_n^2(t_+)$, and
\[
  \frac{h_+}{h_-}=\frac{v_n^2(t_+)}{v_n^2(t_-)}=e^2\qquad\Longrightarrow\qquad e=\sqrt{h_+/h_-}.
\]
Alternatively, one can measure the times between subsequent impacts. After a bounce with upward speed $v_0$,
$x_2(t)=v_0t-\frac12gt^2$, so the flight lasts $t^*=2v_0/g$. The next flight starts with $ev_0$ and lasts
$2ev_0/g$; hence $e=t_2^*/t_1^*$, the ratio of two consecutive flight times.
\end{example}

% =====================================================================
\chapter{Dynamics of Systems of Particles}
% =====================================================================
\lecture{6}{05.10.2026}
We now consider $n$ particles with masses $m_i$ and positions $\vect r_i$, $i=1,\dots,n$, which may interact with
each other. The goal of this chapter is to find balance laws for the system as a whole. The key observation will be
that, thanks to Newton's third axiom, the internal forces cancel from the momentum and angular momentum balances,
so that the motion of the system as a whole is governed by the external forces alone.

\section{Kinematics}

\begin{definition}
The \emph{total mass} and the \emph{centre of mass} of the system are
\begin{equation}
  M=\sum_{i=1}^{n}m_i,\qquad\vect r_\cm=\frac1M\sum_{i=1}^{n}m_i\vect r_i .
\end{equation}
\end{definition}

The particles of a system are frequently connected to each other or to the surroundings. Such connections are
described by constraints, and the most important ones can be written as equations between the positions.

A constraint is \emph{holonomic} if it can be written as an equation between the positions,
\[f(\vect r_1,\dots,\vect r_n,t)=0.\]
A constraint is \emph{holonomic} if it can be written in the form $f(\vect r_1,\dots,\vect r_n,t)=0$.
\end{definition}

A rigid link between particles $i$ and $j$, which cannot change its length, is described by the holonomic constraint
$|\vect r_i(t)-\vect r_j(t)|=\text{const}$, and a rigid system by the set of constraints
$|\vect r_i-\vect r_j|=\text{const}$ for all pairs $i,j$. Each particle contributes $d$ degrees of freedom, and each
independent constraint removes one:

\begin{proposition}
A system of $n$ particles in $d$ dimensions with $r$ independent holonomic constraints has $nd-r$ degrees of freedom.
\end{proposition}

\begin{example}{2.1}{Counting degrees of freedom}
Two particles connected by a rigid link in three dimensions have $2\cdot3-1=5$ DOFs: three for the position of one
of them and two angles for the direction of the link. Two masses suspended from ropes wound around a stepped pulley
with radii $r_0$ and $r_1$ move by $x_0=r_0\varphi$ and $x_1=r_1\varphi$; the whole system has a single DOF, the
pulley angle $\varphi$. As a last example take four particles in the plane forming a quadrilateral of four rigid
links (a diagonal spring is \emph{not} a constraint), one corner pinned and one on a roller. They have $4\cdot2=8$
coordinates; the pin removes two, the roller one, and the links four, leaving $8-7=1$ DOF.
\end{example}

\section{Kinetics}

Each particle $i$ experiences an external force $\vect F_{i,\ext}$ and internal forces $\vect F_{ij}$ exerted by the
other particles $j$. By Newton's third axiom $\vect F_{ij}=-\vect F_{ji}$, and both act along the line connecting the
particles. The net force on particle $i$ is $\vect F_i=\vect F_{i,\ext}+\sum_{j\neq i}\vect F_{ij}$, and when we sum
over all particles the internal forces cancel pairwise:
\begin{equation}\label{eq:sum-forces}
  \sum_{i=1}^{n}\vect F_i=\sum_{i=1}^{n}\vect F_{i,\ext}.
\end{equation}

\subsection{Balance of linear momentum}
Summing the momentum balances $\vect F_i=\frac{\mathrm d}{\mathrm dt}(m_i\vect v_i)$ of the individual particles and
using \eqref{eq:sum-forces} and the definition of the centre of mass,
\[
  \sum_{i=1}^n\vect F_{i,\ext}=\dv{}{t}\sum_{i=1}^nm_i\vect v_i=\dv{}{t}\Big(M\,\frac1M\sum_{i=1}^nm_i\dot{\vect r}_i\Big)
  =\dv{}{t}\big(M\vect v_\cm\big).
\]
\begin{theorem}
The total momentum of a system of particles is $\vect P=M\vect v_\cm$, and it obeys
\begin{equation}
  \sum_{i=1}^n\vect F_{i,\ext}=\dv{}{t}\big(M\vect v_\cm\big),\qquad\text{so that}\quad\sum_{i=1}^n\vect F_{i,\ext}=M\vect a_\cm
  \quad\text{if $M$ is constant}.
\end{equation}
\end{theorem}
In words: the centre of mass moves like a single particle carrying the total mass, acted upon by the total external
force. No matter how complicated the internal interactions are, they cannot move the centre of mass.

\begin{example}{2.4}{Falling dumbbell}
Two masses $m_1$ and $m_2$ connected by a rigid massless rod are thrown into the air. The only external force is
gravity, so $(m_1+m_2)\vect a_\cm=(m_1+m_2)\vect g$, and the centre of mass moves on the parabola
$\vect r_\cm(t)=\frac12\vect gt^2+\vect v_\cm(0)t+\vect r_\cm(0)$ -- regardless of how the dumbbell tumbles. To
learn about the internal force, the rod force, we have to apply the momentum balance to the individual particles:
$m_1\vect a_1=m_1\vect g+\vect F_{12}$ and $m_2\vect a_2=m_2\vect g+\vect F_{21}$ with $\vect F_{21}=-\vect F_{12}$.
\end{example}

\subsection{Work--energy balance}
Summing the work--energy balances $T_i(t_2)-T_i(t_1)=W_{i,12}$ of the individual particles gives the work--energy
balance of the system,
\begin{equation}
  T(t_2)-T(t_1)=W_{12},\qquad T=\frac12\sum_{i=1}^nm_i|\vect v_i|^2,\qquad
  W_{12}=\sum_{i=1}^n\int_{\vect r_i(t_1)}^{\vect r_i(t_2)}\vect F_i\cdot\dd\vect r_i .
\end{equation}
There is an important difference to the momentum balance: \emph{internal forces may do work}. The two forces of a
compressed spring between two particles are equal and opposite, but the particles move differently, and the work
of the two forces does not cancel. Accordingly, $\vect F_i$ here contains the internal forces. As before, conservative
forces can be accounted for by a potential, and for a conservative system $T+V=\text{const}$.

There is, however, an important class of internal forces that do no work.

\begin{proposition}
A rigid link never performs work.
\end{proposition}
\begin{proof}
Differentiating the constraint $|\vect r_i-\vect r_j|^2=\text{const}$ gives
$2(\vect r_i-\vect r_j)\cdot(\vect v_i-\vect v_j)=0$. The link forces act along the link,
$\vect F_{ij}\parallel\vect r_i-\vect r_j$, and $\vect F_{ji}=-\vect F_{ij}$, so their combined work is
\[
  \int\vect F_{ij}\cdot\dd\vect r_i+\int\vect F_{ji}\cdot\dd\vect r_j=\int_{t_1}^{t_2}\vect F_{ij}\cdot(\vect v_i-\vect v_j)\dd t=0 .\qedhere
\]
\end{proof}

\subsection{Balance of angular momentum}
The angular momentum of the system about a point $B$ is the sum of the individual angular momenta,
$\vect H_B=\sum_i\vect r_{BP_i}\times\vect p_i$. Repeating the derivation of \eqref{eq:amb-particle},
\[
  \dot{\vect H}_B=\sum_{i}(\vect v_{P_i}-\vect v_B)\times\vect p_i+\sum_i\vect r_{BP_i}\times\vect F_i
  =-\vect v_B\times\vect P+\sum_i\vect r_{BP_i}\times\vect F_i ,
\]
since $\vect v_{P_i}\times\vect p_i=\vect 0$ and $\sum_i\vect p_i=\vect P$. The internal forces come in pairs that are
equal, opposite and collinear, so their torques cancel, and only the torque $\vect M_B$ of the external forces remains.

\begin{theorem}
For a system of particles, $\vect M_B=\dot{\vect H}_B+\vect v_B\times\vect P$ with $\vect P=M\vect v_\cm$, where
$\vect M_B$ is the net external torque about $B$.
\end{theorem}

\begin{example}{2.5}{Particles on a disk}
Particles $m_i$ are attached at radii $r_i$ to a massless disk which rotates by the angle $\varphi(t)$ about its fixed
centre $B$. Each particle moves on a circle, $\vect v_i=r_i\dot\varphi\,\vect e_{\varphi,i}$, so
\[
  \vect H_B=\sum_{i=1}^nr_i\vect e_{r,i}\times m_ir_i\dot\varphi\,\vect e_{\varphi,i}=\Big(\sum_{i=1}^nm_ir_i^2\Big)\dot\varphi\,\e3 ,
\]
and the angular momentum balance about the fixed point $B$ gives
$M_B^{\ext}=I_B\ddot\varphi$ with $I_B=\sum_im_ir_i^2$. If the external torque is due to gravity, it can be
simplified: $\vect M_B^{\ext}=\sum_i\vect r_{BP_i}\times m_i\vect g=\big(\sum_im_i\vect r_{BP_i}\big)\times\vect g=
M\vect r_{B,\cm}\times\vect g$. In a homogeneous gravitational field, the weight acts as if the total mass were
concentrated at the centre of mass.
\end{example}

\begin{example}{2.6}{A ring of particles}
For $n$ particles of mass $M/n$ evenly distributed on a circle of radius $R$, $I_B=\sum_i(M/n)R^2=MR^2$. In the limit
$n\to\infty$ this is the moment of inertia of a thin ring about its centre.
\end{example}

\lecture{7}{07.10.2026}
\begin{example}{2.8}{Atwood machine}
Two masses $m_1$ (left) and $m_2$ (right) hang from a rope passing over a massless wheel of radius $R$, whose bearing
$B$ is fixed; the wheel is driven by a torque $M_T$. We seek the acceleration and the bearing force.

If $x_1$ is the downward displacement of $m_1$ and $x_2$ the upward displacement of $m_2$, then $x_1=R\varphi=x_2\eqqcolon x$;
the system has one DOF. We cut the rope on both sides of the wheel and call the rope forces $S_1$ and $S_2$. The
momentum balances of the two masses are
\[
  m_1\ddot x=m_1g-S_1,\qquad m_2\ddot x=-m_2g+S_2 .
\]
The wheel is massless, so the forces and torques on it balance: the bearing force is $B_2=S_1+S_2$ (vertically) and
$B_1=0$, and the angular momentum balance about $B$ reads $0=I_B\ddot\varphi=S_1R-S_2R+M_T$. Adding the two
momentum balances and using $S_2-S_1=M_T/R$,
\[
  (m_1+m_2)\ddot x=(m_1-m_2)g+\frac{M_T}{R}\qquad\Longrightarrow\qquad
  \ddot x=\frac{m_1-m_2}{m_1+m_2}\,g+\frac{M_T}{(m_1+m_2)R}. \tag{$\ast$}
\]
The bearing force is
\[
  B_2=S_1+S_2=(m_1+m_2)g+(m_2-m_1)\ddot x .
\]
The first term is the static value; the second appears only because the system moves. In particular, the
static result $B_2=(m_1+m_2)g$ is wrong in dynamics.

The same result follows from the work--energy balance, once we know how much work a torque does. A torque $M_T$ is
equivalent to a couple of forces $F$ applied at distance $R$ on either side of the axis, $M_T=2FR$. When the wheel
turns by $\dd\varphi$ each force moves by $R\dd\varphi$, so the work is $2FR\dd\varphi=M_T\dd\varphi$. Hence
\[
  \frac{m_1+m_2}{2}\dot x^2+m_2gx-m_1gx-[\text{initial value}]=M_T(\varphi-\varphi_0),
\]
and differentiating with respect to time, with $\dot\varphi=\dot x/R$, gives
$(m_1+m_2)\dot x\ddot x+(m_2-m_1)g\dot x=M_T\dot x/R$, i.e.\ again $(\ast)$.

Finally, we may also treat the two masses together with the wheel as one system. Placing the origin at the axle,
the masses hang at horizontal positions $\mp R$ and their heights change by $\mp x$, so
$\vect a_\cm=\frac{m_2-m_1}{m_1+m_2}\ddot x\,\e2$. The momentum balance of the whole system,
$(m_1+m_2)\vect a_\cm=B_1\e1+(B_2-m_1g-m_2g)\e2$, reproduces $B_1=0$ and the expression for $B_2$.
\end{example}

% ---------------------------------------------------------------------
\section{Collisions and variable mass}
\subsection{Particle collisions}
Two smooth spheres collide. The plane tangent to both at the point of contact is the \emph{collision plane}; let
$\e n$ be its normal and $\e t$ a direction in it. Since the spheres are smooth, the contact impulses are normal:
particle 2 receives $\hat F\e n$ and particle 1, by Newton's third axiom, $-\hat F\e n$:
\[
  m_1\vect v_1(t_+)-m_1\vect v_1(t_-)=-\hat F\e n,\qquad m_2\vect v_2(t_+)-m_2\vect v_2(t_-)=\hat F\e n .
\]
As for the wall, the tangential velocity components are unchanged, $v_i^t(t_+)=v_i^t(t_-)$. For the normal components
we again split the impact into compression and restitution:
\begin{align*}
  m_1v_1^n(t_0)-m_1v_1^n(t_-)&=-\hat F_{\text{comp}}, & m_2v_2^n(t_0)-m_2v_2^n(t_-)&=\hat F_{\text{comp}},\\
  m_1v_1^n(t_+)-m_1v_1^n(t_0)&=-\hat F_{\text{rest}}, & m_2v_2^n(t_+)-m_2v_2^n(t_0)&=\hat F_{\text{rest}}.
\end{align*}
These are four equations for six unknowns (everything except the initial velocities). The two missing equations
are the definition of the coefficient of restitution, $e=\hat F_{\text{rest}}/\hat F_{\text{comp}}$, and the
condition that at the instant of maximum compression both bodies have the same normal velocity,
$v_1^n(t_0)=v_2^n(t_0)\eqqcolon v_0$.

To solve, we add the two compression equations; $\hat F_{\text{comp}}$ drops out, and
$v_0=\big(m_1v_1^n(t_-)+m_2v_2^n(t_-)\big)/(m_1+m_2)$ turns out to be the normal velocity of the centre of mass.
Then $\hat F_{\text{comp}}=m_2\big(v_0-v_2^n(t_-)\big)=\frac{m_1m_2}{m_1+m_2}\big(v_1^n(t_-)-v_2^n(t_-)\big)$, the total
normal impulse is $(1+e)\hat F_{\text{comp}}$, and we obtain
\begin{align}
  v_1^n(t_+)&=\frac{m_1-e\,m_2}{m_1+m_2}\,v_1^n(t_-)+\frac{(1+e)\,m_2}{m_1+m_2}\,v_2^n(t_-),\\
  v_2^n(t_+)&=\frac{m_2-e\,m_1}{m_1+m_2}\,v_2^n(t_-)+\frac{(1+e)\,m_1}{m_1+m_2}\,v_1^n(t_-).
\end{align}
Subtracting these two expressions shows that the coefficient of restitution is the ratio of the relative normal
velocities after and before the collision,
\begin{equation}
  e=-\frac{v_2^n(t_+)-v_1^n(t_+)}{v_2^n(t_-)-v_1^n(t_-)} .
\end{equation}
Moreover, since there are no external forces in the normal direction, the momentum is conserved,
$m_1v_1^n(t_+)+m_2v_2^n(t_+)=m_1v_1^n(t_-)+m_2v_2^n(t_-)$.

Three special cases are worth remembering (Examples 2.9 and 2.10). Letting $m_2\to\infty$ with $v_2^n=0$ we recover the
rigid wall, $v_1^n(t_+)=-ev_1^n(t_-)$. In a plastic collision, $e=0$, the two bodies move on together with the velocity
of the centre of mass. And in an elastic collision of equal masses, $e=1$ and $m_1=m_2$, the two particles simply
exchange their normal velocities -- the principle of Newton's cradle.

\subsection{Particles with variable mass}
A body that gains or loses mass -- a rocket burning fuel, a raindrop collecting moisture -- requires care, because the
momentum balance applies to a \emph{fixed} collection of matter. Consider a body of mass $m(t)$ and velocity $\vect v(t)$
which, during the time $\dd t$, absorbs a mass $\dd m$ arriving with the velocity $\vect v_m$. We apply the momentum
balance to the closed system consisting of the body and the mass element. At time $t$ its momentum is
$\vect p(t)=m\vect v+\dd m\,\vect v_m$, and at time $t+\dd t$, using Taylor expansions,
\[
  \vect p(t+\dd t)=\big[m+\dot m\dd t+O(\dd t^2)\big]\big[\vect v+\dot{\vect v}\dd t+O(\dd t^2)\big]
  =m\vect v+\dot m\vect v\dd t+m\dot{\vect v}\dd t+O(\dd t^2).
\]
With $\dd m=\dot m\dd t$ the rate of change of momentum is
\[
  \dot{\vect p}=\lim_{\dd t\to0}\frac{\vect p(t+\dd t)-\vect p(t)}{\dd t}=\dot m\vect v+m\dot{\vect v}-\dot m\vect v_m .
\]
\begin{theorem}[Rocket equation]
For a body of variable mass,
\begin{equation}
  \sum_i\vect F_i=m\vect a+\dot m(\vect v-\vect v_m)=m\vect a-\dot m\,\vect v_{\text{rel}},
  \qquad\text{equivalently}\qquad m\vect a=\sum_i\vect F_i+\dot m\,\vect v_{\text{rel}},
\end{equation}
where $\vect v_{\text{rel}}=\vect v_m-\vect v$ is the velocity of the gained or ejected mass relative to the body, and
$\dot m\,\vect v_{\text{rel}}$ is the \emph{thrust}.
\end{theorem}
Note that the naive formula $\sum\vect F=\frac{\mathrm d}{\mathrm dt}(m\vect v)$ is \emph{not} correct for a body of
variable mass: it would give different results in different inertial frames. Only the relative velocity of the mass
flow matters.

\begin{example}{2.11}{Rocket launch}
A rocket moves vertically upwards, $\vect v=\dot x\e2$, and ejects its exhaust gas downwards with the velocity
$\vect v_{\text{rel}}=-c\e2$ ($c>0$) relative to the rocket. It is subject to gravity and to a drag $F_{\text{drag}}$.
The rocket equation gives
\[
  (-F_{\text{drag}}-mg)\,\e2=m\ddot x\,\e2+\dot mc\,\e2\qquad\Longrightarrow\qquad m\ddot x=-F_{\text{drag}}-mg-\dot mc .
\]
Since the rocket loses mass, $\dot m<0$, the last term $-\dot mc$ is positive; it is the thrust, and the rocket
accelerates upwards only if it exceeds $mg+F_{\text{drag}}$.
\end{example}

% =====================================================================
\chapter{Dynamics of Rigid Bodies}\label{ch:rigid}
% =====================================================================
\lecture{8}{12.10.2026}
A rigid body can be thought of as the limit of a system of particles, $n\to\infty$ with $\Delta m_i\to0$, in which all
mutual distances are fixed: $|\vect r_{AB}|=\text{const}$ for any two points $A,B$ of the body $\mathcal B$. Everything
we derived for systems of particles therefore carries over, with sums replaced by integrals. What is new is that a rigid
body has an orientation, and its rotation must be described and accounted for. In the kinematics this leads to the
angular velocity vector and the transfer formulas; in the kinetics it leads to the moment of inertia tensor.

\section{Kinematics}

With the density $\rho=\mathrm dm/\mathrm dV$, the total mass and the centre of mass of a body are
\begin{equation}
  M=\lim\sum_{i=1}^n\Delta m_i=\int_{\mathcal B}\dd m=\int_{\mathcal B}\rho\dd V,\qquad
  \vect r_\cm=\frac1M\int_{\mathcal B}\vect r\dd m .
\end{equation}
How many degrees of freedom does a rigid body have? In the plane its position is fixed by the two coordinates of one
of its points and one angle; in space by three coordinates and three angles.

\begin{proposition}
An unconstrained rigid body in $d$ dimensions has $d(d+1)/2$ degrees of freedom: $d$ translations and $d(d-1)/2$
rotations (1 DOF in 1D, 3 in 2D, 6 in 3D). Consequently $n$ rigid bodies with $k$ independent holonomic constraints
have $n\,d(d+1)/2-k$ DOFs.
\end{proposition}

\subsection{Velocity and acceleration of points of a rigid body}
For two points $A$ and $B$ of the body we have $\vect r_B=\vect r_A+\vect r_{AB}$, hence
$\vect v_B=\vect v_A+\dot{\vect r}_{AB}$ and $\vect a_B=\vect a_A+\ddot{\vect r}_{AB}$. Since the length of
$\vect r_{AB}$ is constant, the vector can only change its direction; its derivatives must therefore be expressible
through the rotation of the body (and vanish if the body does not rotate).

To make this precise, recall from Mechanics I the angular velocity vector $\vect\omega(t)=\omega(t)\,\e\omega(t)$,
whose direction $\e\omega$ is the axis of rotation; in plane motion $\e\omega=\e3$ is constant and points out of the
plane. If the body rotates about a fixed point $\Pi$, every point moves according to $\vect v=\vect\omega\times\vect r_{\Pi P}$.
Subtracting this relation for two points $A$ and $B$, $\vect v_B-\vect v_A=\vect\omega\times(\vect r_{\Pi B}-\vect r_{\Pi A})=\vect\omega\times\vect r_{AB}$.
The result does not refer to $\Pi$ any more and holds for general rigid motion.

\begin{theorem}[Velocity transfer formula]
For any two points $A$ and $B$ of a rigid body,
\begin{equation}\label{eq:vtransfer}
  \vect v_B=\vect v_A+\vect\omega\times\vect r_{AB}.
\end{equation}
\end{theorem}

\begin{example}{3.1}{Rolling disk}
\begin{center}
\begin{tikzpicture}[scale=1.0,>=Stealth]
  \draw[thick] (0,1) circle (1);
  \fill[pattern=north east lines] (-2,0) rectangle (3,-0.15); \draw (-2,0)--(3,0);
  \fill (0,1) circle (1.5pt) node[below left]{cm};
  \fill (0,0) circle (1.5pt) node[below right]{$A$};
  \fill (0,2) circle (1.5pt) node[above left]{$T$};
  \draw[->,thick] (0,1)--(1,1) node[right]{$v_\cm$};
  \draw[->,thick] (0,2)--(2,2) node[right]{$2v_\cm$};
  \draw[dashed] (0,0)--(2,2);
  \draw[->,thick] (-0.75,1.85) arc (130:60:0.9); \node at (-0.7,2.25){$\omega$};
\end{tikzpicture}
\end{center}
A disk of radius $R$ rolls without slipping. Rolling means that the distance travelled by the centre equals the
arc length unrolled, $x_\cm=R\varphi$, so $v_\cm=R\omega$. For rolling to the right the disk turns clockwise,
$\vect\omega=-\omega\e3$, and the transfer formula gives for the top point $T$ and the contact point $A$
\begin{align*}
  \vect v_T&=R\omega\e1+(-\omega\e3)\times R\e2=2R\omega\e1,\\
  \vect v_A&=R\omega\e1+(-\omega\e3)\times(-R\e2)=\vect 0 .
\end{align*}
The contact point is momentarily at rest, and the velocities along the vertical diameter grow linearly from $0$ to
$2v_\cm$, exactly as if the disk were rotating about $A$.
\end{example}

The last observation is general. In plane motion with $\omega\neq0$ there is always a point $\Pi$ of the body (or of
its rigid extension) with $\vect v_\Pi=\vect 0$, the \emph{instantaneous centre of rotation} (ICR), and the body
momentarily rotates about it: $\vect v_P=\vect\omega\times\vect r_{\Pi P}$ for every point $P$. If the velocity
$\vect v_A$ of one point $A$ is known, the ICR is found at
\begin{equation}
  \vect r_{A\Pi}=\frac1\omega\,\e\omega\times\vect v_A ,
\end{equation}
on the line through $A$ perpendicular to $\vect v_A$, at the distance $|\vect v_A|/\omega$. Indeed, by
\eqref{eq:vtransfer}, $\vect v_\Pi=\vect v_A+\e\omega\times(\e\omega\times\vect v_A)=\vect v_A-\vect v_A=\vect 0$,
since $\e\omega\times(\e\omega\times\vect v)=-\vect v$ for any $\vect v$ in the plane.

A second useful consequence of rigidity is obtained by differentiating $\vect r_{AB}\cdot\vect r_{AB}=\text{const}$:
\[
  2\,\dot{\vect r}_{AB}\cdot\vect r_{AB}=0\quad\Longleftrightarrow\quad(\vect v_B-\vect v_A)\cdot\vect r_{AB}=0 .
\]
\begin{proposition}[Theorem of projected velocities]
For any two points of a rigid body, $\vect v_A\cdot\vect r_{AB}=\vect v_B\cdot\vect r_{AB}$: their velocity
components along the connecting line are equal.
\end{proposition}

\begin{example}{3.2}{Slider--crank mechanism}
A crank $AB$ of length $L$ turns about the fixed pivot $A$ with angular velocity $\vect\omega_2$ and makes the angle
$\alpha$ with the ground. The connecting rod $BC$ of length $2L$, inclined at the angle $\beta$, drives a slider $C$
along the ground, $\vect v_C=-v_C\e1$. Given $\omega_2$, we want $v_C$.

With $\e1$ to the right and $\e2$ upwards the geometry is
$\vect r_{AB}=-L\cos\alpha\,\e1+L\sin\alpha\,\e2$ and $\vect r_{BC}=-2L\cos\beta\,\e1-2L\sin\beta\,\e2$, where
$L\sin\alpha=2L\sin\beta$ because $C$ and $A$ are at the same height. Since $A$ is fixed and
$\vect\omega_2=\omega_2\e3$, the velocity of $B$ is
\[
  \vect v_B=\vect\omega_2\times\vect r_{AB}=-L\omega_2\big(\sin\alpha\,\e1+\cos\alpha\,\e2\big).
\]
The quickest route to $v_C$ is the theorem of projected velocities applied to the rod, $\vect v_B\cdot\vect r_{BC}=\vect v_C\cdot\vect r_{BC}$:
\[
  2L^2\omega_2\big(\sin\alpha\cos\beta+\cos\alpha\sin\beta\big)=2Lv_C\cos\beta
\]
and therefore
\[
  v_C=\omega_2L\big(\sin\alpha+\cos\alpha\tan\beta\big)=\omega_2L\,\frac{\sin(\alpha+\beta)}{\cos\beta}.
\]
Alternatively one can use the transfer formula twice,
$\vect v_C=\vect v_A+\vect\omega_2\times\vect r_{AB}+\vect\omega_1\times\vect r_{BC}=-v_C\e1$, which gives two scalar
equations for $v_C$ and the angular velocity $\vect\omega_1$ of the rod. In terms of the angles,
$\vect\omega_1=\dot\beta\,\e3$ and $\vect\omega_2=-\dot\alpha\,\e3$.
\end{example}

Differentiating the velocity transfer formula with respect to time, and using $\dot{\vect r}_{AB}=\vect\omega\times\vect r_{AB}$,
gives the corresponding relation for accelerations.

\begin{theorem}[Acceleration transfer formula]
\begin{equation}
  \vect a_B=\vect a_A+\dot{\vect\omega}\times\vect r_{AB}+\vect\omega\times\big(\vect\omega\times\vect r_{AB}\big).
\end{equation}
In plane motion the last term simplifies to the centripetal term $-\omega^2\vect r_{AB}$.
\end{theorem}

\begin{example}{3.3}{Rolling disk, accelerations}
With $\vect a_\cm=R\dot\omega\e1$ and $\dot{\vect\omega}=-\dot\omega\e3$,
\begin{align*}
  \vect a_T&=R\dot\omega\e1+(-\dot\omega\e3)\times R\e2-\omega^2R\e2=2R\dot\omega\e1-R\omega^2\e2,\\
  \vect a_A&=R\dot\omega\e1+(-\dot\omega\e3)\times(-R\e2)-\omega^2(-R\e2)=R\omega^2\e2 .
\end{align*}
The contact point, although momentarily at rest, is accelerated upwards towards the centre. The ICR has zero velocity,
but in general \emph{not} zero acceleration -- a fact that has to be kept in mind whenever the ICR is used in kinetics.
\end{example}

\lecture{9}{14.10.2026}
\begin{example}{3.5}{Ladder sliding down a wall}
Two blocks of mass $m$ are connected by a massless rigid rod of length $L$; block $A$ slides on a frictionless
vertical wall and block $B$ on the frictionless floor. Let $\varphi$ be the angle of the rod with the vertical. The system
is released from rest from (almost) vertical. We want the velocities, accelerations and forces.

Two particles in the plane have four coordinates; the wall, the floor and the rod impose three constraints, so there
is a single DOF, $\varphi(t)$. With $\vect r_A=L\cos\varphi\,\e2$ and $\vect r_B=L\sin\varphi\,\e1$,
\[
  \vect v_A=-L\dot\varphi\sin\varphi\,\e2,\qquad\vect v_B=L\dot\varphi\cos\varphi\,\e1 .
\]
The system is conservative, so energy conservation (potential zero at floor level; only $A$ changes height) gives
\[
  mgL=mgL\cos\varphi+\frac m2\big(v_A^2+v_B^2\big)=mgL\cos\varphi+\frac m2L^2\dot\varphi^2,
\]
hence
\[
  \dot\varphi^2=\frac{2g}{L}\,(1-\cos\varphi),
\]
and differentiating in time, $\ddot\varphi=(g/L)\sin\varphi$.

Differentiating the velocities (or, equivalently, using the acceleration transfer formula with
$\vect r_{AB}=L\sin\varphi\,\e1-L\cos\varphi\,\e2$) gives the downward acceleration of $A$ and the horizontal
acceleration of $B$:
\begin{align*}
  a_A&=L\ddot\varphi\sin\varphi+L\dot\varphi^2\cos\varphi=g\sin^2\varphi+2g(1-\cos\varphi)\cos\varphi\\
     &=g\big(1+2\cos\varphi-3\cos^2\varphi\big),\\
  a_B&=L\ddot\varphi\cos\varphi-L\dot\varphi^2\sin\varphi=g\sin\varphi\cos\varphi-2g(1-\cos\varphi)\sin\varphi\\
     &=g\sin\varphi\,(3\cos\varphi-2).
\end{align*}
To find the forces, let $S$ be the rod force (positive in tension), $H$ the force of the wall on $A$ and $V$ the force of the
floor on $B$. The horizontal momentum balance of $B$, in which $V$ does not appear, and the vertical one of $A$ read
\[
  ma_B=-S\sin\varphi,\qquad ma_A=mg+S\cos\varphi\qquad\Longrightarrow\qquad S=mg\,(2-3\cos\varphi),
\]
and the remaining balances give $V=mg-S\cos\varphi$ and $H=-S\sin\varphi$.

The result tells an interesting story. As long as $\cos\varphi>2/3$, $S<0$: the rod is in compression and pushes $A$
against the wall, $H>0$. At $\cos\varphi=2/3$, i.e.\ $\varphi\approx48.2^\circ$, the wall force vanishes; beyond
that point the wall would have to pull on the block, which it cannot do. Block $A$ separates from the wall and the
above solution ceases to be valid.
\end{example}

\subsection{Rigid body kinematics in three dimensions}
The velocity transfer formula holds unchanged in three dimensions, but now the axis of rotation $\e\omega(t)$ may
change its direction; at each instant it defines the \emph{instantaneous axis of rotation} (IAR). In contrast to the
ICR in the plane, the points on the IAR need not be at rest: they may move along the axis,
$\vect v_{\text{IAR}}=v_{\text{IAR}}\,\e\omega$, so that the general instantaneous motion is a screw. If, however, two points
of a body are at rest, the IAR passes through both of them.

\begin{example}{3.6}{Rolling cone}
A cone has its apex $O$ fixed on a vertical axis at height $R$; its axis $\e r$ is horizontal with length $L$, and its base
circle of radius $R$ rolls on the ground, touching it at $A$. The cone precesses about the vertical with the angle
$\theta$ and spins about its own axis with the angle $\varphi$. Given $\dot\theta$, we seek $\vect\omega$.

For any point $C$ of the cone, $\vect v_C=\vect v_O+\vect\omega\times\vect r_{OC}=\vect\omega\times\vect r_{OC}$, since
$O$ is fixed. Rolling means that the contact point is at rest, so $\vect\omega\times\vect r_{OA}=\vect 0$, and as both
vectors are non-zero, $\vect\omega\parallel\vect r_{OA}$: with two points at rest, the line $OA$ \emph{is} the IAR.

The angular velocity is the sum of the precession about the vertical and the spin about the cone axis,
$\vect\omega=\dot\theta\,\e3-\dot\varphi\,\e r$. Since $\vect r_{OA}=L\e r-R\e3$, parallelism requires
$\dot\theta/\dot\varphi=R/L=\tan\alpha$, where $\alpha$ is the half-angle of the cone; this is the rolling condition
$R\dot\varphi=L\dot\theta$. Therefore
\[
  \vect\omega=\dot\theta\,\e3-\frac LR\,\dot\theta\,\e r .
\]
The example illustrates that angular velocities are additive: the total rotation is the vector sum of the individual
rotations.
\end{example}

% =====================================================================
\section{Kinetics}

\subsection{Balance of linear momentum}
Passing to the limit in the momentum of a system of particles, the momentum of a rigid body is
\[
  \vect P=\lim\sum_{i=1}^n\vect v_i\Delta m_i=\int_{\mathcal B}\vect v\dd m=M\vect v_\cm ,
\]
where the last equality follows by differentiating $\vect r_\cm=\frac1M\int_{\mathcal B}\vect r\dd m$. The momentum
balance of the system carries over without change.

\begin{theorem}
$\dfrac{\mathrm d}{\mathrm dt}\vect P=\sum_i\vect F_i^{\ext}$, and for constant mass $M\vect a_\cm=\sum_i\vect F_i^{\ext}$.
\end{theorem}
As far as its centre of mass is concerned, a rigid body is a ``super-particle''; under gravity alone, for instance,
$\vect a_\cm=\vect g$. The momentum balance, however, says nothing about the rotation. For that we need the angular
momentum balance, which for a rigid body is a genuinely independent law.

\subsection{Balance of angular momentum}
\lecture{10}{19.10.2026}
The angular momentum balance for systems, $\vect M_B=\dot{\vect H}_B+\vect v_B\times\vect P$ with $\vect P=M\vect v_\cm$,
also carries over. The work lies in computing the angular momentum of a rigid body about a point $B$ of the body,
\[
  \vect H_B=\lim\sum_{i=1}^n\vect r_{BP_i}\times\Delta m_i\vect v_i=\int_{\mathcal B}(\vect r-\vect r_B)\times\vect v\dd m .
\]
The velocity of each point is given by the transfer formula, $\vect v=\vect v_B+\vect\omega\times(\vect r-\vect r_B)$. Inserting
it, the integral splits into two parts. The first is
\[
  \int_{\mathcal B}(\vect r-\vect r_B)\dd m\times\vect v_B=\Big[\int_{\mathcal B}\vect r\dd m-\vect r_B\int_{\mathcal B}\dd m\Big]\times\vect v_B
  =M(\vect r_\cm-\vect r_B)\times\vect v_B .
\]
For the second, write $\bar{\vect r}=\vect r-\vect r_B$ and use $\vect a\times(\vect b\times\vect c)=\vect b(\vect a\cdot\vect c)-\vect c(\vect a\cdot\vect b)$:
\[
  \bar{\vect r}\times(\vect\omega\times\bar{\vect r})=(\bar{\vect r}\cdot\bar{\vect r})\,\vect\omega-\bar{\vect r}\,(\bar{\vect r}\cdot\vect\omega)
  =\big[(\bar{\vect r}\cdot\bar{\vect r})\Id-\bar{\vect r}\bar{\vect r}^{\mathsf T}\big]\vect\omega ,
\]
where $\Id$ is the unit tensor and $\bar{\vect r}\bar{\vect r}^{\mathsf T}$ the dyadic product. The angular velocity is the
same for all points of the body and can be taken out of the integral. What remains depends only on the distribution of
mass in the body, and deserves a name.

\begin{definition}
The \emph{(mass) moment of inertia tensor} of the body $\mathcal B$ with respect to the point $B$ is
\begin{equation}
  \tens I_B=\int_{\mathcal B}\big[(\bar{\vect r}\cdot\bar{\vect r})\Id-\bar{\vect r}\bar{\vect r}^{\mathsf T}\big]\dd m,
  \qquad\bar{\vect r}=\vect r-\vect r_B .
\end{equation}
\end{definition}

\begin{theorem}[Angular momentum balance for a rigid body]
The angular momentum of a rigid body about a point $B$ of the body is
$\vect H_B=M(\vect r_\cm-\vect r_B)\times\vect v_B+\tens I_B\vect\omega$, and
\begin{equation}
  \vect M_B=\dv{}{t}\Big[M(\vect r_\cm-\vect r_B)\times\vect v_B+\tens I_B\vect\omega\Big]+\vect v_B\times\vect P .
\end{equation}
\end{theorem}

In components, with $\bar{\vect r}=(\bar x_1,\bar x_2,\bar x_3)$, we have
$(\bar{\vect r}\cdot\bar{\vect r})=\bar x_1^2+\bar x_2^2+\bar x_3^2$ and $(\bar{\vect r}\bar{\vect r}^{\mathsf T})_{ij}=\bar x_i\bar x_j$,
so that
\begin{equation}
  [\tens I_B]=\int_{\mathcal B}\begin{bmatrix}
    \bar x_2^2+\bar x_3^2 & -\bar x_1\bar x_2 & -\bar x_1\bar x_3\\
    -\bar x_1\bar x_2 & \bar x_1^2+\bar x_3^2 & -\bar x_2\bar x_3\\
    -\bar x_1\bar x_3 & -\bar x_2\bar x_3 & \bar x_1^2+\bar x_2^2
  \end{bmatrix}\rho\dd V .
\end{equation}
The diagonal entries are the moments of inertia about the coordinate axes -- integrals of the squared distance from
the respective axis -- and the off-diagonal entries are the products of inertia.

The general form of the balance is rarely needed, because a suitable choice of $B$ makes it much simpler. If $B$ is the
centre of mass, the first term of $\vect H_B$ vanishes, and so does $\vect v_B\times\vect P=\vect v_\cm\times M\vect v_\cm$;
if $B$ is a fixed point, $\vect v_B=\vect 0$ and the same terms vanish. In both cases
\begin{equation}\label{eq:amb-simple}
  \vect M_B=\dv{}{t}\big(\tens I_B\vect\omega\big),
\end{equation}
and one should always choose $B$ in one of these two ways. If, in addition, $\tens I_B$ is constant, then
$\vect M_B=\tens I_B\dot{\vect\omega}$, the three-dimensional counterpart of the planar law $M_B=I_B\dot\omega$; and if
there is no torque, the angular momentum $\tens I_B\vect\omega$ is conserved.

Why must the moment of inertia be a tensor, rather than a number? If the coordinate axes happen to be chosen such that
$\tens I_\cm$ is diagonal, the balance decouples into $M_{\cm,i}=I_{\cm,i}\dot\omega_i$, $i=1,2,3$, and each axis
has its own moment of inertia. In general, however, the angular momentum $\tens I\vect\omega$ is not parallel to the
angular velocity: the resistance of a body against rotation depends on the axis, and the linear map that captures this
dependence is a second-order tensor.

\begin{example}{3.7}{Moment of inertia of a cylinder}
Consider a homogeneous cylinder of radius $R$ and length $L$, mass $M=\rho\pi R^2L$, with its axis along $x_3$ and the
origin at its centre of mass. In cylindrical coordinates $x_1=r\cos\varphi$, $x_2=r\sin\varphi$, $x_3$, with
$\dd V=r\dd r\dd\varphi\dd x_3$,
\[
  [\tens I_\cm]=\int_0^R\!\!\int_{-L/2}^{L/2}\!\!\int_0^{2\pi}
  \begin{bmatrix}r^2\sin^2\varphi+x_3^2&-r^2\sin\varphi\cos\varphi&-rx_3\cos\varphi\\
  &r^2\cos^2\varphi+x_3^2&-rx_3\sin\varphi\\ \text{sym.}&&r^2\end{bmatrix}\rho\,r\dd\varphi\dd x_3\dd r .
\]
All off-diagonal entries vanish, because $\int_0^{2\pi}\sin\varphi\cos\varphi\dd\varphi=\int_0^{2\pi}\cos\varphi\dd\varphi=0$.
For the axial moment, $I_{33}=\rho\cdot\frac{R^4}{4}\cdot2\pi\cdot L=M\frac{R^2}{2}$. The transverse moment $I_{11}$ consists
of two parts: $\rho\frac{R^4}{4}\cdot\pi\cdot L=M\frac{R^2}{4}$ from the term $r^2\sin^2\varphi$, and
$\rho\frac{R^2}{2}\cdot2\pi\cdot\frac{L^3}{12}=M\frac{L^2}{12}$ from the term $x_3^2$. Hence
\[
  [\tens I_\cm]=\begin{bmatrix}\frac{MR^2}{4}+\frac{ML^2}{12}&0&0\\0&\frac{MR^2}{4}+\frac{ML^2}{12}&0\\0&0&\frac{MR^2}{2}\end{bmatrix}.
\]
\end{example}

Two limiting cases of this result are used constantly in plane problems. A disk or cylinder rotating about its own
axis has $I_\cm=\frac12MR^2$, and a slender rod or bar ($R\to0$) rotating about an axis through its centre perpendicular
to it has $I_\cm=\frac1{12}ML^2$. In these planar applications we have tacitly assumed $\dot I_\cm=0$. This is justified:
for rotation about $x_3$ the balance reads $M_{\cm,3}=\frac{\mathrm d}{\mathrm dt}(I_{\cm,33}\omega_3)$, and
$I_{\cm,33}=\int_{\mathcal B}(\bar x_1^2+\bar x_2^2)\rho\dd V=\int_{\mathcal B}\bar r^2\rho\dd V$ only involves the distances
from the axis, which do not change in a plane rotation. Hence $M_\cm=I_\cm\dot\omega$.

\begin{example}{}{Pulling up a weight}
A pulley of mass $M$ and radius $R$ on a fixed axle ($I_\cm=\frac12MR^2$) is driven by a torque $T$ and lifts a weight
$m$ by a rope. With the rope force $S$ and the height $x=R\varphi$ of the weight,
\[
  m\ddot x=S-mg,\qquad I_\cm\ddot\varphi=T-SR .
\]
Eliminating $S=T/R-\frac12M\ddot x$ from the first equation gives
\[
  \ddot x=\frac{T/R-mg}{m+M/2}.
\]
The pulley behaves like an additional mass $M/2$ that has to be accelerated together with the weight.
\end{example}

% ---------------------------------------------------------------------
\subsection{Properties of the moment of inertia tensor}
\lecture{11}{21.10.2026}
Computing moments of inertia by direct integration quickly becomes laborious. Fortunately the tensor has a number of
properties that reduce most practical calculations to tabulated values and a little algebra.

\paragraph{Additivity.}
Since an integral over a union of regions is the sum of the integrals, the moment of inertia of a body composed of parts
$\mathcal B_1$ and $\mathcal B_2$ is the sum of the moments of inertia of the parts -- provided all of them refer to the
\emph{same} point $B$:
\begin{equation}
  \tens I_B=\tens I_{B,1}+\tens I_{B,2}.
\end{equation}
Holes can be treated by subtraction. A hollow cylinder, for example, is a full cylinder of radius $R_o$ minus one of
radius $R_i$, and by Example 3.7
\[
  I_{11}=I_{22}=\frac{M_o}{12}\big(L^2+3R_o^2\big)-\frac{M_i}{12}\big(L^2+3R_i^2\big),\qquad
  I_{33}=\frac{M_oR_o^2}{2}-\frac{M_iR_i^2}{2},
\]
with $M_{o,i}=\rho\pi R_{o,i}^2L$. For a thin ring with $R_{o,i}=R\pm t/2$ and $t\to0$ the axial moment becomes
$I_{33}=\frac{\rho\pi L}{2}(R_o^2+R_i^2)(R_o^2-R_i^2)\to(\rho\,2\pi RtL)\,R^2=MR^2$, in agreement with Example 2.6.

\paragraph{Symmetry.}
For a homogeneous body, symmetry forces many entries of the tensor to vanish.
\begin{proposition}
If the $x_i$-axis is an axis of rotational symmetry of the body, or if it is normal to a plane of symmetry, then all
off-diagonal components with index $i$ vanish (the reference point must lie on the axis or in the plane).
\end{proposition}
Consider, for instance, a cone with its apex at $B$ and its axis along $x_1$. Since $x_1$ is an axis of rotational
symmetry, $I_{12}=I_{13}=0$; since $x_2$ and $x_3$ are normal to planes of symmetry, also $I_{23}=0$. The reason is
easily seen in the integral $I_{12}=-\int_{\mathcal B}x_1x_2\rho\dd V$: every slice $x_1=\text{const}$ is a circle centred
on the axis, and for each of its points with coordinate $x_2$ there is one with $-x_2$, so the contributions cancel.
The same argument applies to $I_{23}=-\int x_2x_3\rho\dd V$.

\paragraph{Positive definiteness and principal axes.}
For an arbitrary unit vector $\vect n$,
\[
  \vect n\cdot\tens I_B\vect n=\int_{\mathcal B}\big[(\vect r\cdot\vect r)(\vect n\cdot\vect n)-(\vect n\cdot\vect r)(\vect r\cdot\vect n)\big]\rho\dd V
  =\int_{\mathcal B}\big[|\vect r|^2-(\vect r\cdot\vect n)^2\big]\rho\dd V>0 ,
\]
because the projection of $\vect r$ onto $\vect n$ is never longer than $\vect r$ itself; in fact the integrand is the
squared distance of the point from the axis through $B$ in the direction $\vect n$. The tensor is therefore symmetric and
positive definite, and linear algebra tells us that it can be diagonalised with real, positive eigenvalues.

\begin{theorem}[Principal axes]
For every body $\mathcal B$ and every point $B$ there is a frame $C$ in which
\[
  [\tens I_B]_C=\begin{bmatrix}I_1&&\\&I_2&\\&&I_3\end{bmatrix},\qquad I_1,I_2,I_3>0 .
\]
The axes of this frame, the eigenvectors of $\tens I_B$, are the \emph{principal axes} of the body, and the
$I_k$ its principal moments of inertia. Axes of rotational symmetry and axes normal to symmetry planes are always
principal axes.
\end{theorem}
In practice one should always try to work in principal axes; with respect to the centre of mass the balance then
takes the simple form $[\vect M_\cm]=(I_1\dot\omega_1,\,I_2\dot\omega_2,\,I_3\dot\omega_3)^{\mathsf T}$.

For plane motion only the component $I_{B,33}=\int_{\mathcal B}(x_1^2+x_2^2)\rho\dd V=\int_{\mathcal B}r^2\rho\dd V$ is
needed, and the balance reads $M_B=I_{B,33}\dot\omega$. For a long slender rod with its pivot $B$ at one end, for
example, $x_2$ is negligible across the thin cross-section $A$, and
\[
  I_B=\int_0^Lx_1^2\rho A\dd x_1=\frac13\rho AL^3=\frac13ML^2 .
\]

\paragraph{The parallel axis theorem.}
Tables give moments of inertia with respect to the centre of mass. To transfer them to another point we use the
following result, known in German as the \emph{Satz von Steiner}.

\begin{theorem}[Parallel axis theorem]
Let $\Delta\vect x=\vect r_B-\vect r_\cm$. Then
\begin{equation}
  \tens I_B=\tens I_\cm+M\big[(\Delta\vect x\cdot\Delta\vect x)\Id-\Delta\vect x\,\Delta\vect x^{\mathsf T}\big],
\end{equation}
that is, in components,
\[
  [\tens I_B]=[\tens I_\cm]+M\begin{bmatrix}
    \Delta x_2^2+\Delta x_3^2 & -\Delta x_1\Delta x_2 & -\Delta x_1\Delta x_3\\
    -\Delta x_1\Delta x_2 & \Delta x_1^2+\Delta x_3^2 & -\Delta x_2\Delta x_3\\
    -\Delta x_1\Delta x_3 & -\Delta x_2\Delta x_3 & \Delta x_1^2+\Delta x_2^2
  \end{bmatrix}.
\]
In plane motion, $I_B=I_\cm+M|\Delta\vect x|^2$.
\end{theorem}
\begin{proof}
Write $\vect r-\vect r_B=\vect s-\Delta\vect x$ with $\vect s=\vect r-\vect r_\cm$ and insert this into the definition of
$\tens I_B$. Expanding the products, the terms quadratic in $\vect s$ give $\tens I_\cm$, the terms quadratic in
$\Delta\vect x$ give the second term (with $\int\dd m=M$), and the mixed terms are linear in $\vect s$ and vanish, since
$\int_{\mathcal B}\vect s\dd m=\vect 0$ by the definition of the centre of mass.
\end{proof}
The proof shows that the theorem relates an arbitrary point to the \emph{centre of mass}; to move between two arbitrary
points one has to pass through the centre of mass. It also shows that, among all parallel axes, the one through the
centre of mass has the smallest moment of inertia.

Two examples: for the slender rod, $I_\cm=\frac1{12}ML^2$ and $|\Delta\vect x|=L/2$ give
$I_B=\frac1{12}ML^2+\frac14ML^2=\frac13ML^2$, in agreement with the direct integration above. For a wheel of radius $R$,
$I_\cm=\frac12MR^2$, and about its contact point $A$ with the ground $I_A=I_\cm+MR^2=\frac32MR^2$. For a rolling
wheel the angular momentum balance may be used in either of the forms $M_\cm=I_\cm\dot\omega$ and $M_A=I_A\dot\omega$
(see the remark after Example 3.10).

\paragraph{Relation to area moments of inertia.}
For a thin plate of thickness $t$ lying in the $x_1x_2$-plane, the mass moments of inertia are proportional to the
area moments of inertia known from beam theory. The polar moment is $I_p=\int r^2\dd A=I_{33}/(\rho t)$; since $x_3$ is
negligible across the plate, $I_{11}=\int(x_2^2+x_3^2)\rho\dd V\approx\rho t\int x_2^2\dd A=\rho t\,I_x$ and similarly
$I_{22}=\rho t\,I_y$; and the product of inertia is $I_{12}=-\rho t\,I_{xy}$.

% ---------------------------------------------------------------------
\lecture{12}{26.10.2026}
The following example shows how the balance laws are combined in a system of several rigid bodies.

\begin{example}{3.10}{Pulley system}
Pulley 1 (mass $m$, radius $r$, $I_\cm=\frac12mr^2$) turns on a fixed axle at the ceiling and is driven by a force $F$
acting tangentially on its rim. A rope runs from pulley 1 down under the free pulley 2 (mass $M$, radius $R$) and back up
to a fixed point on the ceiling. A weight of mass $m_w$ hangs from the centre of pulley 2. We call the rope forces $S_1$
(between the pulleys), $S_2$ (to the ceiling) and $S_3$ (between pulley 2 and the weight).

The kinematics is governed by one observation: the rope segment leading to the ceiling does not move, so the point
where it leaves pulley 2 is the ICR of that pulley. If pulley 2 and the weight move up by $x$, then $x=R\theta$, and the
opposite side of pulley 2, at distance $2R$ from the ICR, moves twice as fast. The rope, and with it the rim of pulley 1,
therefore moves by $2x$: $r\varphi=2R\theta=2x$. The system has one DOF.

For the weight, the momentum balance gives
\[
  m_w\ddot x=S_3-m_wg. \tag{I}
\]
For pulley 1, the angular momentum balance about its fixed centre, with $\ddot\varphi=2\ddot x/r$, gives
\[
  \tfrac12mr^2\cdot\frac{2\ddot x}{r}=(F-S_1)\,r\qquad\Longrightarrow\qquad m\ddot x=F-S_1 . \tag{II}
\]
For pulley 2 we take moments about its ICR, which eliminates the unknown $S_2$. With $I_{\text{ICR}}=\frac32MR^2$ and
$\ddot\theta=\ddot x/R$,
\[
  \tfrac32MR^2\,\ddot\theta=S_1\cdot2R-(S_3+Mg)\,R\qquad\Longrightarrow\qquad\tfrac32M\ddot x=2S_1-Mg-S_3 . \tag{III}
\]
The combination $2\cdot\text{(II)}+\text{(III)}+\text{(I)}$ eliminates $S_1$ and $S_3$:
\[
  \Big(2m+\tfrac32M+m_w\Big)\ddot x=2F-Mg-m_wg\qquad\Longrightarrow\qquad\ddot x=\frac{2F-(M+m_w)g}{2m+\frac32M+m_w}.
\]
\end{example}

\begin{remark}
The angular momentum balance about a moving point in general contains additional terms, and the simple form
$M=I\dot\omega$ about the ICR is not justified in general. For a rolling disk (or a pulley rolling on its rope) it is:
the contact point is momentarily at rest and its acceleration points towards the centre of mass, and with these two
properties the additional terms vanish. The form $M_A=I_A\dot\omega$ may therefore be used for rolling wheels, but it
should not be applied blindly to other moving points.
\end{remark}

\subsection{The angular momentum transfer formula}
The balance $\vect M_B=\dot{\vect H}_B+\vect v_B\times\vect P$ holds for any point $B$, which need not even belong to the
body. It is often convenient to relate angular momenta about different points. Writing $\vect r_B=\vect r_\cm+\vect r_{\cm,B}$,
\[
  \vect H_B=\int_{\mathcal B}(\vect r-\vect r_B)\times\vect v\dd m
  =\int_{\mathcal B}(\vect r-\vect r_\cm)\times\vect v\dd m-\vect r_{\cm,B}\times\int_{\mathcal B}\vect v\dd m
  =\vect H_\cm+\vect P\times\vect r_{\cm,B}.
\]
Writing the same relation for a second point $A$ and subtracting, with $\vect r_{\cm,B}-\vect r_{\cm,A}=\vect r_{AB}$, we
obtain:
\begin{theorem}[Angular momentum transfer formula]
\begin{equation}
  \vect H_B=\vect H_A+\vect P\times\vect r_{AB}.
\end{equation}
\end{theorem}
It has exactly the structure of the moment transfer formula of statics, $\vect M_B=\vect M_A+\vect F\times\vect r_{AB}$,
with the force replaced by the momentum.

\begin{example}{I}{Disk rolling without slipping}
A disk of mass $M$ and radius $R$ on rough ground (static friction coefficient $\mu_0$) is pulled by a horizontal force $F$
applied at its centre. The ground exerts the normal force $N$ and a friction force $F_R$ at the contact point $A$. The
momentum balance in the horizontal direction is $Ma_\cm=F-F_R$. For the angular momentum balance we use the contact
point, about which neither $N$ nor $F_R$ has a torque: $I_A\dot\omega=FR$ with $I_A=\frac32MR^2$. Thus
\[
  \dot\omega=\frac23\frac F{MR},\qquad a_\cm=R\dot\omega=\frac23\frac FM,\qquad F_R=F-Ma_\cm=\frac13F .
\]
Without friction the disk would slide without rotating, since $F$ has no torque about the centre of mass; it is the
friction force that makes the disk turn. Because the contact point is the ICR and does not slide, this is \emph{static}
friction, and rolling without slipping is only possible as long as $F_R\le\mu_0N=\mu_0Mg$, i.e.\ $F\le3\mu_0Mg$.
\end{example}

\begin{example}{II}{Spinning disk placed on the ground}
A disk spinning with $\omega(0)=\omega_0$ but with $v_\cm(0)=0$ is placed on the ground (friction coefficient $\mu$).
At first it slips, so dynamic friction of magnitude $F_R=\mu Mg$ acts at the contact point, in the direction in which it
pushes the centre forward. The momentum balance and the angular momentum balance about the centre of mass give
\[
  Ma_\cm=\mu Mg\;\Rightarrow\;v_\cm(t)=\mu gt,\qquad
  \tfrac12MR^2\dot\omega=-\mu MgR\;\Rightarrow\;\omega(t)=\omega_0-\frac{2\mu g}{R}\,t .
\]
The disk speeds up while its rotation slows down, until the rolling condition $v_\cm=R\omega$ is reached at the time
$t_f$; from then on the contact point no longer slides, friction becomes static (and in fact zero, since no horizontal
force is needed to keep a disk rolling), and the disk rolls on uniformly. From $\mu gt_f=R\omega_0-2\mu gt_f$,
\[
  t_f=\frac{R\omega_0}{3\mu g},\qquad v_\cm(t_f)=\frac13R\omega_0,\qquad\omega(t_f)=\frac13\omega_0 .
\]
There is an elegant way to obtain the final state without knowing the friction force. Take the angular momentum balance
about the contact point $A$, which moves along with the disk. Its velocity is parallel to $\vect P$, so
$\vect v_A\times\vect P=\vect 0$; and all external forces either act at $A$ or (gravity and $N$) cancel, so
$\vect M_A=\vect 0$. Hence $\vect H_A$ is conserved. (Note that $\vect M_A=\vect 0$ does \emph{not} imply
$\dot\omega=0$: while the disk slips the contact point is not the ICR, and $\vect H_A\neq I_A\vect\omega$.) By the
transfer formula,
\begin{align*}
  \vect H_A&=\vect H_\cm+\vect P\times\vect r_{\cm,A}=-\tfrac12MR^2\omega\,\e3+Mv_\cm\e1\times(-R\e2)\\
  &=-\big(\tfrac12MR^2\omega+MRv_\cm\big)\e3 ,
\end{align*}
so $\frac R2\omega+v_\cm$ is constant, and comparing the initial state with the rolling state $v_\cm=R\omega$ gives
$\frac R2\omega_0=\frac32v_\cm$, i.e.\ $v_\cm=\frac13R\omega_0$ as before.
\end{example}

% ---------------------------------------------------------------------
\subsection{Work--energy balance}
\lecture{13}{28.10.2026}
To formulate the work--energy balance we first need the kinetic energy of a rigid body. Choose a reference point $C$ of
the body and insert the transfer formula $\vect v=\vect v_C+\vect\omega\times(\vect r-\vect r_C)$ into
$T=\frac12\int_{\mathcal B}|\vect v|^2\dd m$. Expanding the square gives three terms:
\[
  T=\frac12|\vect v_C|^2\int_{\mathcal B}\dd m+\vect v_C\cdot\Big(\vect\omega\times\int_{\mathcal B}(\vect r-\vect r_C)\dd m\Big)
  +\frac12\int_{\mathcal B}\big|\vect\omega\times(\vect r-\vect r_C)\big|^2\dd m .
\]
The first integral is $M$. The second is $M(\vect r_\cm-\vect r_C)$, so the middle term vanishes if $C$ is the centre of
mass, and also if $\vect v_C=\vect 0$, i.e.\ if $C$ is the ICR (or lies on the IAR). For the last term we use
$(\vect\omega\times\vect a)\cdot(\vect\omega\times\vect a)=\big[\vect a\times(\vect\omega\times\vect a)\big]\cdot\vect\omega$,
which by the computation of the previous section yields $\frac12(\tens I_C\vect\omega)\cdot\vect\omega$.

\begin{theorem}
If $C$ is the centre of mass or a point with $\vect v_C=\vect 0$, the kinetic energy of a rigid body is
\begin{equation}
  T=\frac12M|\vect v_C|^2+\frac12\,\vect\omega\cdot\tens I_C\vect\omega ,
\end{equation}
the sum of a translational and a rotational part. In plane motion, $T=\frac12Mv_C^2+\frac12I_C\omega^2$.
\end{theorem}

The work--energy balance of the system of particles then becomes
\begin{equation}
  T(t_2)-T(t_1)=W_{12}=\sum_i\int_{\vect r_i(t_1)}^{\vect r_i(t_2)}\vect F_i^{\ext}\cdot\dd\vect r_i+\int_{t_1}^{t_2}\vect M_T\cdot\vect\omega\dd t ,
\end{equation}
where $\vect r_i$ is the point of application of the force $\vect F_i$, and external torques $\vect M_T$ contribute
work as well, as we saw in Example 2.8. Internal forces do no work, since the body is rigid. As before, a conservative
system satisfies $T+V=\text{const}$.

\begin{example}{3.11}{Disk rolling up a hill}
A disk of mass $M$ and radius $R$ rolls without slipping with $v_0=R\omega_0$ onto a slope. How high does it get?

The normal force and the friction force do no work: they act at the contact point, which is the ICR and therefore at
rest. (Static friction on a stationary point does no work.) The system is thus conservative. The kinetic energy can be
computed with respect to the centre of mass,
\[
  T=\frac M2v_\cm^2+\frac12\cdot\frac{MR^2}{2}\Big(\frac{v_\cm}R\Big)^2=\frac34Mv_\cm^2 ,
\]
or, equivalently, with respect to the ICR, using Steiner's theorem:
$T=\frac12\big(\frac12MR^2+MR^2\big)\omega^2=\frac34Mv_\cm^2$. Energy conservation between the bottom and the highest
point gives
\[
  \frac34Mv_0^2=Mgh_{\max}\qquad\Longrightarrow\qquad h_{\max}=\frac34\,\frac{v_0^2}g .
\]
A particle thrown upwards with the same speed would only reach $\frac12v_0^2/g$; a considerable part of the kinetic
energy of the disk is stored in its rotation.
\end{example}

\begin{example}{3.12}{Planetary gear}
A sun gear (radius $r_s$, moment of inertia $I_{\cm,s}$) driven by a torque $M_T$ meshes with three planet gears of radius
$r_p$, which in turn mesh with a fixed outer ring. We seek the angular acceleration $\dot\omega_s$ of the sun gear; gravity
plays no role because of the symmetry of the arrangement.

The contact forces between the gears act with equal magnitude and opposite direction on two points that move together,
so they do no net work; the forces at the fixed ring act at the ICRs of the planets and do no work either. The only work
is done by the driving torque. The kinetic energy is
\[
  T=\frac12I_{\cm,s}\,\omega_s^2+3\cdot\frac12I_{\text{ICR},p}\,\omega_p^2,\qquad I_{\text{ICR},p}=I_{\cm,p}+m_pr_p^2 ,
\]
where we have used that each planet momentarily rotates about its contact point with the ring. The work--energy
balance $T(t)-T(t_0)=\int M_T\dd\varphi_s=M_T(\varphi_s-\varphi_{s,0})$, differentiated with respect to time, reads
\[
  I_{\cm,s}\,\omega_s\dot\omega_s+3I_{\text{ICR},p}\,\omega_p\dot\omega_p=M_T\,\omega_s .
\]
It remains to relate $\omega_p$ to $\omega_s$. The point of contact between sun and planet moves with the speed
$\omega_sr_s$; for the planet this point is at the distance $2r_p$ from its ICR, so $\omega_sr_s=2r_p\omega_p$. Inserting
$\omega_p=\omega_sr_s/(2r_p)$ and cancelling $\omega_s$,
\[
  \dot\omega_s=\frac{M_T}{I_{\cm,s}+3I_{\text{ICR},p}\,\big(r_s/(2r_p)\big)^2}.
\]
The planets act on the sun shaft like an additional, ``reduced'' moment of inertia.
\end{example}

\begin{example}{3.13}{A note on gears}
Gear 1 (radius $r_1$, moment of inertia $I_1$) is driven by a torque $M_1$ and drives gear 2 (radius $r_2$, moment of
inertia $I_2$), which is loaded by a torque $M_2$; $F_r$ is the force between the teeth. Since the gears roll on each other
without slipping, $r_1\omega_1=r_2\omega_2$. The angular momentum balances about the fixed axes are
\[
  I_1\dot\omega_1=M_1-F_rr_1,\qquad I_2\dot\omega_2=F_rr_2-M_2 .
\]
Eliminating $F_r$,
\[
  M_2=M_1\frac{r_2}{r_1}-\Big(I_1\dot\omega_1\frac{r_2}{r_1}+I_2\dot\omega_2\Big),
\]
where the bracket represents the losses due to the inertia of the gears. For massless gears it vanishes, and
$M_1/M_2=r_1/r_2$, which together with the kinematic relation gives $M_1\omega_1=M_2\omega_2$: the power is
transmitted unchanged, and a gear simply trades torque for angular velocity.
\end{example}

\end{document}
